\documentclass[11pt,a4paper]{article}

%----- ENHANCED TYPOGRAPHY -----
\usepackage[utf8]{inputenc}
\usepackage[T1]{fontenc}
\usepackage{lmodern}        % clean vector font
\usepackage{microtype}      % better justification & kerning
\usepackage{palatino} 
\usepackage{braket}   
\usepackage{mathtools}  % in the preamble
\usepackage{bbm}

%----- PAGE LAYOUT -----
\usepackage{geometry}
\geometry{top=1in, bottom=1in, left=1in, right=1in}
\usepackage{setspace}
\onehalfspacing  % 1.5 line spacing

%----- FANCY HEADERS & FOOTERS -----
\usepackage{fancyhdr}
\pagestyle{fancy}
\fancyhf{}
% page number outside, header text inside

\fancyhead[LO]{\small \rightmark}
\fancyhead[RO]{\small \leftmark}
\renewcommand{\headrulewidth}{0.4pt}
\renewcommand{\footrulewidth}{0pt}

% make sections feed into \leftmark/\rightmark
\renewcommand{\sectionmark}[1]{\markboth{#1}{}}
\renewcommand{\subsectionmark}[1]{\markright{#1}}

%----- SECTION NUMBERING & TOC DEPTH -----
\setcounter{secnumdepth}{3}  % number down to \subsubsection
\setcounter{tocdepth}{2}     % show ToC down to \subsection

%----- AMS MATH & THEOREM STYLES -----
\usepackage{amsmath,amssymb,mathtools}
\usepackage{amsthm}

% definitions, examples, remarks upright
\theoremstyle{definition}
\newtheorem{definition}{Definition}[section]
\newtheorem{example}[definition]{Example}
\newtheorem{remark}[definition]{Remark}

% theorems, lemmas, corollaries italic
\theoremstyle{plain}
\newtheorem{theorem}[definition]{Theorem}
\newtheorem{lemma}[definition]{Lemma}
\newtheorem{proposition}[definition]{Proposition}
\newtheorem{corollary}[definition]{Corollary}

% unnumbered proof environment
\theoremstyle{remark}

%----- OTHER PACKAGES -----
\usepackage{graphicx}
\usepackage{tikz}
\usetikzlibrary{arrows.meta,positioning}
\usetikzlibrary{calc, matrix, decorations.pathreplacing, positioning}
\usepackage{tikz-cd}
\usepackage{hyperref}
\hypersetup{colorlinks,
linkcolor=blue, citecolor=purple, urlcolor=teal}
\usepackage{enumitem}
\setlist[itemize]{nosep, left=1.5em}
\usepackage{booktabs}
\usepackage{listings}
\lstset{
basicstyle=\ttfamily\small,
numbers=left,
numbersep=5pt,
frame=single,
breaklines=true
}
\usepackage{xcolor}
\definecolor{shade}{HTML}{F5F5F5}
\usepackage{float}
%----- CUSTOM MACROS -----
\newcommand{\F}{\mathbb{F}}
\newcommand{\code}[1]{\texttt{#1}}
\newcommand{\dist}[2]{d\bigl(#1,#2\bigr)}
\newcommand{\R}{\mathbb{R}}
\newcommand{\Z}{\mathbb{Z}}
\newcommand{\N}{\mathbb{N}} 
\newcommand{\Q}{\mathbb{Q}} 
\newcommand{\C}{\mathbb{C}}
\renewcommand{\set}[1]{\left\{ #1 \right\}}
\newcommand{\angles}[1]{\langle #1 \rangle}
\newcommand{\abs}[1]{\lvert #1 \rvert}
\newcommand{\norm}[1]{\lVert #1 \rVert}
\newcommand{\1}{\mathbbm{1}}
\newcommand{\Span}{\operatorname{Span}}

% \usepackage{mathtools} 
% \DeclarePairedDelimiter{\angles}{\langle}{\rangle} 
% \DeclarePairedDelimiter{\braces}{\left\{}{\right\}} 
% \DeclarePairedDelimiter{\abs}{\lvert}{\rvert} 
% \DeclarePairedDelimiter{\norm}{\lVert}{\rVert}

%----- TITLE METADATA -----
\title{\LARGE\bfseries Functional Analysis}
\author{Georges Khater \\ \small American University of Beirut, Math 305}
\date{\today}

%===============================================
\begin{document}
\maketitle
\tableofcontents
\bigskip

\section{Metric Spaces.} 

\subsection{Topological Spaces}

\begin{definition}
    Given a set $X$, a topology $\tau$ on $X$ is a subset of $\mathcal{P}(X)$ s.t. 
    \begin{enumerate}
        \item $\varnothing, X \in \tau$
        \item Closed under union: $E_\alpha \in \tau \ \forall \alpha \in I \implies \bigcup_{\alpha \in I} E_\alpha \in \tau$
        \item Closed under finite intersection: $E_1, E_2, \cdots, E_n \in \tau \implies \bigcap_{i=1}^n E_i \in \tau$
    \end{enumerate}
    We say that $(X, \tau)$ is a topological space.
\end{definition}

\begin{definition}
    Let $(X, \tau)$ a topological space and $A \subseteq X$
    \begin{itemize}
        \item $A$ open $\iff$ $A \in \tau$. 
        \item $A$ is closed $\iff X \setminus A$ is open.
    \end{itemize}
\end{definition}

\begin{definition}
    $(X, \tau)$ topological space and $x \in X$. A neighborhood of $x$ is an open set $U_x \ni x$. 
\end{definition}

\begin{definition}
    $(X, \tau)$ topological space, ,and $A \subseteq X$. Then $x$ is an interior point of $A$ $\iff \exists U_x$ an open neighborhood of $x$ s.t. 
    $U_x \subseteq X$. 

    We denote by $\mathring{A}$ the set of interior points of $A$. Notice $\mathring{A} \subseteq A$.

    Now $x$ is a limit point of $A$ $\iff \forall U_x$ a neighborhood of $x$, we have that $U_x \setminus \set{x} \cap A \ne \varnothing$. $A'$ denotes the set of all 
    limit points of $A$ and $\overline{A} := A \cap A'$ denotes the closure of $A$.
\end{definition}

\begin{proposition}
    \begin{enumerate}
        \item $\mathring{A}$ is an open set. In fact, $A$ is open $\iff A = \mathring{A}$.
        \item $\overline{A}$ is closed. In fact, $A$ is closed $\iff A' \subseteq A \iff \overline{A} = A$.
    \end{enumerate}

    Consequence, $A$ is open iff $\forall x \in A, \ \exists U_x \ni x$ an open neighborhood s.t. $U_x \subseteq A$.
\end{proposition}

\begin{proof}
    Easy.
\end{proof}

\begin{definition}
    Let $(X, \tau_x), \ (Y, \tau_y)$ topological spaces, and $f \colon X \to Y$ and $x_0 \in X$. Then $f$ is continuous at $x_0 \iff \forall V$ neighborhood of $f(x_0)$ there exists a neighborhood $U$ of $x_0$ s.t. 
    $f(U) \subseteq V$.

    $f$ is said to be continuous on $X$ iff $f$ is continuous at every $x_0 \in X$.
\end{definition}

\begin{proposition}
    $f \colon X \to Y$ continuous iff $f^{-1}(G)$ is open in $X$ for every $G$ open in $Y$. 
\end{proposition}

\begin{definition}
    $(X, \tau)$ topological spaces, and $x_n$ a sequence in $X$. We say that $x_n$ converges to $x_0 \in X$ iff for every $U_{x_0}$ neighborhood of $x_0$, $\exists N \in \N$
    s.t. $x_n \in U_{x_0}$ $\forall n \ge N$. 
\end{definition}

\begin{example}
    $x_n = \frac{1}{n}$. 
    \begin{itemize}
        \item In $(\R, \tau_{\text{trivial}})$, this sequence converges to every point in $\R$.
        \item In $(\R, \tau_{\text{disc}})$, this sequence doesn't converge to any point (we say that it diverges).
    \end{itemize}
\end{example}

\begin{proposition}
    $(X, \tau)$ topological space. $A \subseteq X$, $x_n$ sequence in $A$ that converges to $x_0$. Then $x_0 \in \overline{A}$. 
\end{proposition}

\begin{proposition}
    $f \colon (X, \tau_x) \to (Y, \tau_y)$ and $x_0 \in X$ with $f$ continuous at $x_0$. If $x_n$ converges to $x_0$, then $f(x_n)$ converges to $f(x_0)$.
\end{proposition}

\begin{definition}
    $(X, \tau)$ topological space, $A \subseteq X$. We say that $A$ is dense in $X$ iff $\overline{A} = X$. We say that $X$ is separable iff it has a countable dense subset.
\end{definition}

\begin{example}
    $(\R^n, \abs{\cdot})$ is separable. $(\R, \tau_{\text{disc}})$ is not separable. $(X, \tau_{\text{trivial}})$ is always separable.
\end{example}

\begin{definition}
    $(X, \tau)$ is a Hausdorff space iff $\forall x, y \in X \ x \ne y$ $\exists U_x, U_y$ neighborhoods of $x, y$ respectively s.t. 
    $U_x \cap U_y = \varnothing$.
\end{definition}

\begin{proposition}
    The limit of a convergent sequence in a Hausdorff space is unique.
\end{proposition}

\begin{proof}
    Assume $x_n$ converges to $x, y$ with $x \ne y$, since Hausdorff we have $U_x, U_y$ neighborhoods with $U_x \cap U_y = \varnothing$. Now for $n$ big enough 
    $x_n \in U_x$ for $n \ge N$ then $x_n \not\in U_y$ for these $n \ge N$, contradiction.
\end{proof}

\begin{definition}
    $X$ a set, then a metric on $X$ is a function $d \colon X \times X \to \R$ s.t. 
    \begin{enumerate}
        \item $d(x, y) \ge 0, \ \forall x, y \in X$ with equality iff $x = y$. 
        \item $d(x, y) = d(y,x) \ \forall x, y \in X$. 
        \item $d(x, z) \le d(x,y) + d(y, z) \ \forall x,y,z \in X$
    \end{enumerate} 
    We say that $(X, d)$ is a metric space.
\end{definition}

\begin{definition}
    $x, \in X$ and $r \ge 0$ then 
    $$B(x, r) = \set{y \in X \colon d(x, y) < r}$$ 
    Now define 
    $$\tau_d \set{U \in \tau_d \iff \forall x \in U \exists r \ge 0 s.t. B(x, r) \subseteq U}$$
\end{definition}

\begin{proposition}
    the $B(x, r)$. 
    $\tau_d$ is a topology on $X$, called the metric topology. Then $(X, \tau_d)$ is always Hausdorff and $\tau_d$ is the smallest topology containing all 
\end{proposition}

\begin{proposition}
    Take $(X, d)$ metric space and $x_n$ converges to $x_0$ iff $\forall \varepsilon > 0, \ \exists N \in \N$ s.t. $d(x_n, x_0) < \varepsilon \ \forall n \ge N$.
\end{proposition}

\begin{proposition}
    $(X, d)$ metric space and $E \subseteq X$, then $x \in \overline{E}$ iff $\exists x_n \in E$ a convergent sequence with $x_n \to x$ 
\end{proposition}

\begin{proof}
    \begin{itemize}
        \item[$(\impliedby)$] Any topological space. 
        \item[$(\implies$)] Take $x \in \overline{E} = E \cup E'$, if $x \in E$ take $x_n = x$. Now if $x \in E'$, then for every $n \in \N, \ \exists x_n \in B^*(x, 1/n) \cap E$ then 
        $d(x_n, x) < 1/n$, by Archimedean property we get that $x_n \to x$.
    \end{itemize}
\end{proof}

\begin{proposition}
    $(X, d)$ metric space with $E \subseteq X$, then $E$ is closed iff for every convergent sequence $x_n \in E$ the limit as $n \to \infty$ of $x_n \in E$.
\end{proposition}

\begin{definition}
    $(X, d)$ metric space, then $E \subseteq X$ is bounded iff $E \subseteq B(x_0, R)$ for some $x_0 \in x$ and $R > 0$. In particular, a sequence $x_n$ is bounded 
    iff $\exists M$ s.t. $\abs{x_n} \le M \ \forall n \in \N$.
\end{definition}

\begin{example}[Some examples of some separable or not separable metric spaces]
    \mbox{}\\
    \begin{itemize}
        \item $(\R, \abs{\cdot})$ separable with countable dense set $\Q^n$.
        
        \item The space $C([a,b])$ of real valued continuous functions with $d_\infty(f, g) = \max_{[a, b]} \abs{f(x) - g(x)}$. By Stone-Weierstrass, for every $f \in C([a,b])$ there exists a sequence of polynomials 
        $P_n$ s.t. $\lim_{n \to \infty} P_n = f$ uniformly. Let's call $P(\R)$ the space of polynomials and fix $\varepsilon > 0$ and $f \in C([a,b])$. Then $\exists P \in P(\R)$ s.t. 
        $d_\infty (P, f) < \varepsilon/2$ with 
        $$p(x) = c_k x^k + c_{k-1} x^{k-1} + \cdots + c_1 x + c_0$$
        Now for each $c_i \in q_i \in \Q$ s.t. $\abs{c_i - q_i} < \frac{\varepsilon}{2M}$ where $M \ge \sum_{i=0}^k \abs{x}^i, \forall x \in [a,b]$. Now let 
        $$q(x) = q_k x^k + \cdots + q_1 x + q_0$$
        So we get that 
        \begin{align*}
            \abs{p(x) - q(x)} &\le \sum_{i=0}^k \sum_{q_i - c_i} \abs{x^i} \\ 
            &\le \frac{\varepsilon}{2M} \sum_{i=0}^k \abs{x}^i
        \end{align*}
        So we get that 
        $$d_\infty(q, f) < \varepsilon$$
        Hence $C([a, b])$ is separable.

        \item Let $\ell^\infty (\R)$ be the space of bounded sequences. Now define 
        $$x = (x_1, x_2, \cdots), \quad y = (y_1, y_2, y_3, \cdots)$$ 
        and 
        $$d_\infty (x,y) = \sup_{i \in \N} \abs{x_i - y_i}$$
        Assume it's separable, let $M$ be the countable dense subset. Let $A = \set{0,1}^\N$ be the space of all sequences with entries $0$ or $1$, clearly $A$ is uncountable. But if $x, y \in A$  with $x \ne y$ then 
        $d_\infty (x, y) = 1$. For every $a \in A$, construct $B(a, 1/2)$, let $m_a$ be an element of $B(a, 1/2) \cap M$ (uses AOC and density of $m_a$). Clearly if $a \ne b$, we have 
        $B(a, 1/2) \cap B(b, 1/2) = \varnothing$, then $a \to m_a$ is injective, hence $M$ is uncountable. Contradiction.
    \end{itemize}    
\end{example}

\begin{definition}
    Let $(X, d)$ a metric space, given a sequence $x_n$ in $X$, we say that $x_n$ is Cauchy iff $\forall \varepsilon > 0, \ \exists N$ s.t. 
    $$d(x_n, x_m) < \varepsilon, \ \forall n, m \ge N$$  
\end{definition}

\begin{proposition}
    Every Cauchy sequence is bounded, every convergent sequence is Cauchy.
\end{proposition}

\begin{definition}
    $(X,d)$ is complete iff every Cauchy Sequence converges.
\end{definition}

\begin{example}
    $(\R^n, \abs{\cdot})$ is complete.
\end{example}

Some properties of metric spaces.
\begin{enumerate}[label = \roman*)]
    \item if $x_n$ converges, then $x_n$ is bounded. 
    
    \item Let $C$ be the space of convergent sequences, then $C \subseteq \ell^\infty$. We want to show that $C$ is closed in $\ell^\infty$. 
    Let $x^1 = (x_{1,1}, x_{1,2}, \cdots), x^2 = (x_{2, 1}, x_{2,2}, \cdots), x^3 = (x_{3, 1}, x_{3,2}, \cdots), \cdots$ be a sequence in $C$ that converges to some sequence $x = (x_1, x_2, \cdots) \in \ell^\infty$. 
    Let $\varepsilon > 0$, since $x^n$ converges to $x$ in $\ell^\infty$, $\exists N$ s.t $d_\infty(x^n, x) < \varepsilon/3 \ \forall n \ge N$. In particular for $n = N$ we have 
    $$\abs{x_{N,i} - x_i} < \varepsilon/3, \quad \forall i = 1,2,3,\cdots$$
    Now $x^N$ is Cauchy, then $\exists i_0$ s.t.
    $$\abs{x_{N,i} - x_{N,j}} < \varepsilon/3, \quad \forall i,j \ge i_0$$
    Hence 
    \begin{align*}
        \abs{x_i - x_j} &\le \abs{x_i - x_{N,i}} + \abs{x_{N,i} - x_{N,j}} + \abs{x_{N,j} - x_j} \\
        &< \varepsilon
    \end{align*}
    Then the sequence $x$ is Cauchy, since $\R$ is complete we have that $x \in C$. So $C$ is closed. 

    %TODO: exercises.
    \item Is $\ell^\infty(\R)$ complete. Let $x^1 = (x_{1,1}, x_{2,1}, \cdots), x^2 = (x_{2,1}, x_{2,2}, \cdots), \cdots$ be a Cauchy sequence in $\ell^\infty(\R)$. We define 
    $$x_i := (x_{i}^1, x_{i}^2, x_{i}^3, \cdots)$$
    Clearly we have that for each $i \in \N$, 
    $$d_\infty (x^k, x^l) \ge \abs{x_{k,i} - x_{l, i}}$$
    Hence it must be the case that all the $x_i$'s are Cauchy, and hence converge to some $x_{i}^*$ (since $\R$ is complete). Define 
    $$x^* := (x_1^*, x_2^*, \cdots)$$
    Now by Cauchy we know that $\exists N \in \N$ s.t. 
    $$d_\infty(x_n, x_m) < \varepsilon/2, \quad \forall n,m \ge N$$
    Therefore $\forall i \in \N$ we have that 
    $$\abs{x^n_i - x_i^m} < \varepsilon/2$$
    But we know that $(x^m_i)_i \to x^*_i$, by continuity of $\abs{\cdot}$, we have 
    $$\abs{x^n_i - x_i^*} \le \varepsilon/2, \quad \forall i \in \N, \ \forall n \ge N$$
    Note that $N$ is fixed over all $i$'s hence 
    $$d_\infty (x^n, x^*) = \sup_{i \in \N} \abs{x^n_i - x_i} \le \varepsilon/2 < \varepsilon$$
    Hence $\ell^\infty(\R)$ is complete.

    Note that $x^*$ is bounded because each $x^n$ is bounded in $\ell^\infty$, therefore $\abs{x_{n,i}} \le M, \ \forall n, i$, taking $n \to \infty$ we get that $\abs{x_i} \le M, \ \forall i$.

    \item Is $C([a, b])$ complete. This is just uniformly Cauchy iff uniformly convergent and the uniform limit of continuous functions is continuous.
    
    \item Is $P([a,b])$ closed in $C([a,b])$. No, this is because $P([a,b])$ is dense with nonempty complement. So if its complement was open, it should contain a nonempty open set, contradicting density of 
    the polynomials.
\end{enumerate}

\begin{proposition}
    Let $(X,d_x), \ (Y, d_y)$ metric spaces and $f \colon X \to Y$ and $x_0 \in X$ then $f$ is continuous at $x_0$ iff $\forall \varepsilon > 0, \ \exists \delta > 0$ s.t. 
    $$d_y(f(x), f(x_0)) < \varepsilon, \quad x \text{ s.t. } d_x(x, x_0) < \delta$$
    iff for every sequence $x_n \to x_0$ we have $f(x_n) \to f(x_0)$. 
\end{proposition}

\begin{proposition}
    $(X,d)$ metric space, $E \subseteq X$. If $E$ is complete, then $E$ is closed.
\end{proposition}

\begin{proof}
    Let $x_n$ be a sequence in $E$ that converges to some $x \in X$ , but $x_n$ is Cauchy in $E$ so it converges to $x \in E$ (by uniqueness of the limit).
\end{proof}

\begin{example}
    Note that the converse is false, $(\Q, \abs{\cdot})$ is closed in itself but is not Cauchy.
\end{example}

\begin{proposition}
    Every closed subset of a complete metric space is complete.
\end{proposition}

\begin{proof}
    Take $(X,d)$ complete and $E \subseteq X$ closed. Take a Cauchy sequence $x_n \in E$, since $E \subseteq X$ it converges to some $x \in X$. Now $E$ is closed so we must have $x\in E$. 
\end{proof}

\begin{example}
    $C \subseteq \ell^\infty(\R)$ is complete. 
\end{example}

\begin{remark}
    $(X,d)$ metric space: 
    $$\abs{d(x,y) - d(x',y')} \le d(x,x') + d(y,y')$$

    This is because 
    \begin{align*}
        d(x,y) &\le d(x,x') + d(x', y') + d(y',y) \\
        d(x,y) - d(x',y') &\le d(x,x') + d(y,y') 
    \end{align*}
    Similarly, 
    $$d(x',y) - d(x,y) \le d(x',x) + d(y',y)$$
    So we get the desired inequality.
\end{remark}

\begin{corollary}
    $(X,d)$ metric space and $x_n, y_n$ Cauchy sequences, then $d(x_n, y_n)$ converge in $\R$.
\end{corollary}

\begin{proof}
    Take $d_n = d(x_n, y_n) \in \R$, we show that $d_n$ is Cauchy. Now 
    \begin{align*}
        \abs{d_n - d_m} &= \abs{d(x_n, y_n) - d(x_m, y_m)} \\
        &\le d(x_n, x_m) + d(y_n, y_m) 
    \end{align*}
    Now since $x_n, y_n$ are Cauchy then $d_n$ is Cauchy in $\R$. Which by completeness gives that $d_n$ converges.
\end{proof}

\begin{lemma}
    $(X,d)$ metric space, and $M$ a dense subset of $X$ s.t. every Cauchy sequence in $M$ converges in $X$. Then $X$ is complete.
\end{lemma}

\begin{proof}
    Let $x_n$ be a Cauchy sequence in $X$, for every $n \in \N, \ \exists y_n \in M$ s.t. 
    $$d(x_n, y_n) < \frac{1}{n}$$
    Let $\varepsilon > 0$, now 
    \begin{align*}
        d(y_n, y_m) &\le d(y_n, x_n) + d(x_n, x_m) + d(x_m, y_m) \\
        &< \frac{1}{n} + d(x_n, x_m) + \frac{1}{m}
    \end{align*}
    Now $\exists N$ s.t. $n,m \ge N$ we have 
    $$d(y_n, y_m) < \varepsilon$$
    This is because
    $$1/m, 1/n < \varepsilon/3, \quad d(x_n, x_m) < \varepsilon/3$$
    Now $y_n \in M$ is Cauchy, so $y_n$ converges to some $y_0 \in X$. Now let $\varepsilon$
    \begin{align*}
        d(x_n, y_0) &\le d(x_n, y_n) + d(y_n, y_0) \\
        &< \frac{1}{n} + d(y_n, y_0)
        &< \varepsilon
    \end{align*}
    for $n \ge N$ for some large enough index $N$. Then $x_n \to y_0$ so $X$ is complete.
\end{proof}

\begin{definition}
    $\varphi \colon (X, d_X) \to (Y,d_y)$ is an isometry iff 
    $$d_Y(\varphi(x), \varphi(y)) = d_X(x,y)$$ 
    
    We say that $A$ is isometric to $B$ iff there exists a bijective isometry $\varphi \colon A \to B$.
\end{definition}


\begin{theorem}[Completeness theorem]
    $(X,d)$ metric space, then there exists a complete metric space $(\hat{X}, \hat{d})$ and a dense subset $M$ of $\hat{X}$ s.t. $M$ is isometric to $X$.  

    In fact, the completion is "unique" (up to isometry), we call it the completion of $X$.
\end{theorem}

\begin{proof}
    Let $C_X$ be the space of Cauchy sequences in $X$. We define the relation $\sim$ on $C_X$ as follows 
    $$x_n \sim y_n \iff \lim_{n \to \infty} d(x_n, y_n) = 0$$
    Not that the RHS exists by the above. Clearly $\sim$ is an equivalence relation. Let $\hat{X} = C_X / \sim$ (the set of equivalence classes), i.e. 
    $$\hat{X} = \set{[x_n] \colon x_n \in C_X}$$ 
    We define $\hat{d} ([x_n], [y_n]) = \lim_{n \to \infty} d(x_n, y_n)$ (this exists by the corollary above). Note that this is well-defined, indeed let $x_n \sim x_n'$ and $y_n \sim y_n'$ 
    so we have that 
    \begin{align*}
        \abs{d(x_n, y_n)  - d(x_n', y_n')} \le d(x_n, x_n') + d(y_n, y_n')  \to 0
    \end{align*}
    Therefore 
    $$\lim_{n \to \infty} d(x_n, y_n) = \lim_{n \to \infty} d(x_n', y_n')$$
    We now show that $\hat{d}$ is a metric (clear since $d$ is a metric): 
    \begin{itemize}
        \item Clearly $\hat{d} ([x_n], [y_n]) \ge 0$ and equal to zero iff $\lim_{n \to \infty} d(x_n, y_n) = 0 \iff x_n \sim y_n \iff [x_n] = [y_n]$. 
        \item Symmetry follows from Symmetry of $d$. 
        \item Triangle inequality follows from triangle inequality on $d$. 
    \end{itemize} 
    Define $\varphi \colon X \to \hat{X}$ as follows, 
    $$\varphi(x) = [(x,x,x,\cdots)], \quad (x,x,x,\cdots) \in C_x$$
    This is clearly an isometry, we denote by $M := \varphi(X) = \set{[(x,x,x,\cdots) \colon x \in X]}$. We now show that $\overline{M} = \hat{X}$. 
    Consider an element $[x_n] = [(x_1, x_2, \cdots)] \in \hat{X}$. Let $\varepsilon > 0$, $\exists N \in \N$ s.t. $\forall m,n \ge N$ 
    $$d(x_n, x_m) < \varepsilon /2 $$
    Consider the sequence $P_n = (x_N, x_N, \cdots)$, clearly 
    $$d(x_n, P_n) = d(x_n, x_N) < \varepsilon/2, \quad \forall n \ge N$$
    Then 
    $$\hat{d} ([x_n], [P_n]) = \lim_{n \to \infty} d(x_n, P_n) \le \varepsilon/2 < \varepsilon$$
    With 
    $$[P_n] = [(x_N, x_N, \cdots)] \in M$$ 
    Therefore $\overline{M} = \hat{X}$. It remains to show that $\hat{X}$ is complete, we use the lemma. Take a Cauchy sequence in $M$, 
    $$P_{x_1} = [(x_1, x_1, \cdots)], \ P_{x_2} = [(x_2, x_2, \cdots)], \ \cdots, \ P_{x_k} = [(x_k, x_k, x_k, \cdots)]$$
    Construct the sequence $p = (x_1, x_2, x_3, \cdots)$, note that 
    $$d(x_n, x_m) = \hat{d} ([(x_n)], [(x_m)]) = \hat{d} (P_{x_n}, P_{x_m})$$
    But $P_{x_n}$ is Cauchy in $\hat{X}$, hence $p$ is Cauchy in $X$, so $P := [p] \in \hat{X}$. Take $\varepsilon > 0$, then $\exists N \in \N$ s.t. 
    $$d(x_n, x_k) \le \varepsilon/2, \quad \forall n, k \ge N$$
    So we get that for $k \ge N$
    \begin{align*}
        \hat{d} (P, P_{x_k}) &= \lim_{n \to \infty} d(x_n, x_k) \\
        &\le \varepsilon/2 \\
        &< \varepsilon
    \end{align*}
    Hence 
    $$\lim_{k \to\infty} P_{x_k} = P \in \hat{X}$$
\end{proof}

\subsection{Compactness}

\begin{definition}
    $(X, \tau)$ topological spaces and $K \subseteq X$. Then $K$ is compact iff every open cover $\mathcal{G} = \set{G_\alpha}_{\alpha \in I}$ of $K$ has a finite subcover.
\end{definition}

\begin{example}
    $[0,1)$ is not compact in $(\R, \abs{\cdot})$, $\mathcal{G} = \set{(-1, 1-1/n) \colon n \in \N}$ has no finite subcover.

    Finite sets are clearly compact. 
\end{example}

\begin{proposition}
    $(X, \tau)$ topological space, then every closed subset of a compact set is compact. 
\end{proposition}

\begin{proof}
    Let $K$ be compact, and $E \subseteq K$ is closed. Let $\mathcal{G} = \set{G_\alpha}_{\alpha \in I}$ open cover of $E$, so 
    $$E \subseteq \bigcup_{\alpha \in I} G_\alpha$$
    But $E^c$ is open with 
    $$K \subseteq \left(\bigcup_{\alpha \in I} G_\alpha \right)\cup E^c$$
    Hence $\set{G_\alpha} \cup \set{E^c}$ is an open cover of $K$, so it has a finite subcover 
    $$K \subseteq G_{\alpha_1} \cup \cdots \cup G_{\alpha_N} \cup E^c$$
    Then clearly 
    $$E = E \cap K \subseteq G_{\alpha_1} \cup \cdots \cup G_{\alpha_N}$$
    so $E$ is compact. 
\end{proof}

\begin{proposition}
    Let $(X, \tau)$ be a Hausdorff space, then $K$ compact $\implies$ $K$ is closed. 
\end{proposition}

\begin{proof}
    Take $q \in K^c$, for every $p \in K$ by Hausdorff $\exists U_p, V_{p}^q$ open neighborhoods of $p, q$ resp. s.t. 
    $U_p \cap V_p^q = \varnothing$. Take $\mathcal{G} = \set{U_p}_{p \in K}$ an open cover of $K$, so it has a finite subcover 
    $$K \subseteq U_{p_1} \cup \cdots \cup U_{p_N}$$
    Now define 
    $$V^q = V_{p_1}^q \cap \cdots \cap V_{p_N}^q$$
    open set with 
    $$V^q \cap K = \varnothing$$
    So $K$ is closed. 
\end{proof}

\begin{proposition}
    $(X,d)$ metric space, $K \subseteq X$ is compact $\implies$ $K$ is closed and bounded.
\end{proposition}

\begin{proof}
    Ex.
\end{proof}

\begin{theorem}[Heine Borel]
    The converse is true for $(\R^n, \abs{\cdot})$ or more generally f.d. vector spaces.
\end{theorem}

\begin{example}
    $\ell^\infty(\R) \ d_\infty(x, y) = \sup_{x_i - y_i}$ and $A = \set{0,1}^\N$. Note that $A$ is bounded, since if $x_0 = (0,0,0,0,\cdots)$, notice that 
    $d_\infty(x, x_0 ) \le 1$, so $A$ is bounded. Moreover, every convergent sequence in $A$ is eventually a constant and so belongs to $A$ so $A$ is closed. 
    We could also say that $A^c$ is open, take $x \not\in A$, then $\exists x_i \ne 0, \ x_i \ne 1$. So $\exists \varepsilon>0$ s.t. forall $y \in (x_i - \varepsilon, x_i + \varepsilon)$ $y_0 \ne 0, \ y \ne 1$. 
    take $z \in B(x, \varepsilon)$ so 
    $$\abs{z_i - x_i} \le d_\infty(z,x) < \varepsilon$$
    i.e. $z_i \ne 0$ and $z_i \ne 1$ so $z \in A^c$. 

    Now $A$ is not compact. Take $\mathcal{G} = \set{B(x,1/2)}_{x \in A}$, now $B(x,1/2) \cap A = \set{x}$ but $A$ is infinite (uncountable even), then $\mathcal{G}$ does not have a finite subcover. 
\end{example}

\begin{definition}
    $(X, d)$ metric space and $K \subseteq X$, then $K$ is totally bounded iff $\forall \varepsilon > 0, \ \exists x_1, x_2, \cdots, x_n \in K$ s.t. $K \subseteq B(x_1, \varepsilon) \cup B(x_2, \varepsilon) \cup \cdots \cup B(x_n, \varepsilon)$. 
\end{definition}

\begin{example}
    $\set{0,1}^\N$ is not totally bounded in $\ell^\infty(\R)$ (although it is bounded).
\end{example}

\begin{proposition}
    $K$ is totally bounded $\implies K$ is bounded.
\end{proposition}

\begin{proof}
    easy.
\end{proof}

\begin{proposition}
    $A$ subset of a totally bounded set is totally bounded. 
\end{proposition}

\begin{proof}
    Let $K$ be totally bounded and $E \subseteq K$, fix $\varepsilon >0$ then we have $x_1, x_2, \cdots, x_n \in K$ s.t. 
    $$K \subseteq B(x_1, \varepsilon/2) \cup \cdots \cup B(x_n, \varepsilon/2)$$
    Note that $E \subseteq K$, so pick $y_1, \cdots, y_k \in B(x_i, \varepsilon/2)$ (in case of empty just skip). Then we have that 
    $$E \subseteq B(x_1, \varepsilon) \cup \cdots \cup B(x_k, \varepsilon)$$  
\end{proof}

\begin{lemma}
    Assume that $x_n$ is a Cauchy sequence in $(X, d)$ with a subsequence that converges to $x_0$, then $x_n \to x_0$.    
\end{lemma}

\begin{proof}
    easy.
\end{proof}

\begin{remark}
    In $\R^n$ (finite dimensional vector spaces), $K$ is totally bounded iff $K$ is bounded. 
\end{remark}

\begin{proof}
    Let $K$ be bounded, then $\exists Q$ a closed cube in $\R^n$ s.t. $K \subseteq Q$. Now $Q$ is compact (closed and bounded), so given $\varepsilon > 0$ it can be covered by finitely many balls of radius $\varepsilon$ and center in $\Q$. 
    So the cube is totally bounded, and hence $K$ is a subset of a totally bounded set, hence is totally bounded. 
\end{proof}

\begin{theorem}
    $(X, d)$ metric space and $K \subseteq Q$, then the following are equivalent: 
    \begin{enumerate}[label = \roman*)]
        \item $K$ is compact 
        \item Every infinite subset $E$ of $K$ has a limit point in $K$. 
        \item Every sequence in $K$ has a convergent subsequence (with its limit in $K$). 
        \item $K$ is complete and totally bounded. 
    \end{enumerate}
\end{theorem}

\begin{proof}
    \begin{itemize}
        \item[$(i) \implies (ii)$] Let $E \subset K$ be infinite s.t. $E$ has no limit points in $K$. Then for each $x \in K$, $x \not\in E'$ so $\exists \varepsilon_x$ s.t. 
        $$B^*(x, \varepsilon_x) \cap E = \varnothing$$
        Now take $\mathcal{G} = \set{B(x, \varepsilon_x)}_{x \in K}$, this is an open cover of $K$ so it has a finite subcover. We get 
        $$K \subseteq B(x_1, \varepsilon_{x_1}) \cup \cdots \cup B(x_N, \varepsilon_{x_N})$$
        But $E \subseteq K$ so 
        $$E = E \cap K \subseteq \bigcup_{i=1}^N  E \cap B(x_1, \varepsilon_{x_i}) \subseteq \set{x_1, \cdots, x_N}$$
        But $E$ is infinite, contradiction.

        \item[$(ii) \implies (iii)$] Let $x_n$ be a sequence in $K$ and define 
        $$E = \set{x_n\colon n \in \N}$$
        If $E$ is finite, by pigeonhole principle $\exists n_1 < n_2 < \cdots$ and $x_0 \in E$ s.t. $x_{n_k} = x_0 \in E$ so we have found a convergent subsequence. If not then by the hypothesis $E$ has a limit point $x_0 \in K$.
        We know that 
        $$B(x_0, 1) \cap E \ne \varnothing$$
        So let $n_1$ be such that $x_{n_1} \in B(x_0, 1) \cap E$. In fact if $x \in E'$ we have that $B(x, \varepsilon) \cap E$ is infinite (since we are in a metric space). Hence for each $k$, 
        $$B^*(x_0, 1/k) \cap E \text{ is infinite}$$
        So we can find some $n_k > n_{k-1}$ with $d(x_{n_k}, x_0) < 1/k$, we have that 
        $$\lim_{k \to \infty} x_{n_k} = x_0$$  

        \item[$(iii) \implies (iv)$] Let $x_n$ be a Cauchy sequence in $K$. But $K$ is sequentially compact, then $x_n$ has a convergent subsequence $x_{n_k} \to x_0 \in K$. 
        Now by the lemma, $x_n$ converges to $x_0$ and so $K$ is complete. Assume $K$ is not totally bounded, then $\exists \varepsilon_0 > 0$ s.t. $K$ cannot be covered by finitely many $\varepsilon_0$-balls with center in $K$. 
        take $x_1 \in K$, then take $x_2 \in K \setminus B(x_1, \varepsilon_0)$, then $x_3 \in K \setminus (B(x_1, \varepsilon_0) \cup B(x_2, \varepsilon_0))$. Similarly take 
        $$x_k \in K \setminus \left(B(x_1, \varepsilon_0) \cup \cdots \cup B(x_{k-1}, \varepsilon_0)\right)$$
        Consider the sequence $\set{x_n}$, then for $n \ne m$ we have 
        $$d(x_n, x_m) \ge \varepsilon_0$$
        So $x_n$ has no convergent subsequence, contradiction.  

        \item[$(iv) \implies (i)$] Assume $K$ is not compact, then $\exists \mathcal{G} = \set{G_\alpha}_{\alpha \in I}$ an open cover of $K$ that does not have a finite subcover. 
        Take $\varepsilon = 1$, so $K$ can be covered by finitely many $1$-balls. So a$\exists x_1 \in K$ and $B(x_1, 1)$ s.t. 
        $K_1 := B(x_1, 1) \cap K$ cannot be covered by a finite subcover of $\mathcal{G}$. Now take $\varepsilon = 1/2$, since $K_1 \subseteq K$ is totally bounded we can similarly find $x_2 \in K_1$ with 
        $K_2 := B(x_2, 1/2) \cap K_1$ cannot be covered by finitely many $G_\alpha$'s. Recursively, we construct 
        $$K_0 \subseteq K_1 \subseteq K_2 \subseteq \cdots$$
        with $K_n := B(x_n, 1/n) \cap K_{n-1}$ with $x_{n} \in K_{n-1}$ and $K_n$ cannot be covered by finitely many elements in $\mathcal{G}$. We show that $x_n$ is Cauchy. 
        Fix $\varepsilon > 0$, take $N$ s.t. $2/N < \varepsilon$, then for $n, m > N$ then 
        $$x_n \in K_{n-1} \subseteq K_N, \quad x_m \in K_{m-1} \subseteq K_N \implies x_n, x_m \in K_N = B(x_{n}, 1/N) \cap K_{N-1}$$
        So 
        \begin{align*}
            d(x_n, x_m) &\le d(x_n, x_N) + d(x_N, x_m) \\
            &< 2/N\\
            &< \varepsilon
        \end{align*}
        So $x_n$ is Cauchy.
    \end{itemize}
    Therefore by completeness, $x_n$ converges to some $x_0 \in K$. But we must have that $\exists \alpha_0 \in I$ s.t. 
    $x_0 \in G_{\alpha_0}$ which is open, so $\exists \delta_0 > 0$ s.t. $B(x_0, \delta_0) \subseteq G_{\alpha_0}$. Since $x_n \to x_0$, then $\exists N$ s.t. 
    $1/N < \delta_0/2$ and $d(x_n, x_0) < \delta_0/2$. So take $y \in B(x_n, 1/N)$ so 
    \begin{align*}
        d(y,x_0) &\le d(y,x_n) + d(x_n, x_0) \\
        &< 1/N + d(x_n, x_0) \\
        &< \delta_0
    \end{align*}
    So 
    $$B(x_n, 1/N) \subseteq B(x_0, \delta_0)$$
    We get that $K_N = B(x_N, 1/N) \cap K_{N-1} \subseteq B(x_0, \delta_0) \subseteq G_{\alpha_0}$. 
\end{proof}

\begin{proposition}
    $f \colon (X, \tau_X) \to (Y, \tau_Y)$ and $K \subseteq X$ compact, then $f(K)$ is compact.
\end{proposition}

\begin{proof}
    easy.
\end{proof}

\subsection{The Arzela-Ascoli theorem} 

\begin{definition}
    Let $X$ be a compact metric space, we denote by $C(X)$ the space of real-valued continuous functions on $X$. We define 
    $$d_\infty(f, g) = \sup_{x \in X} \abs{f(x) - g(x)}$$
\end{definition}

\begin{proposition}
    $C(X)$ is complete and separable.
\end{proposition}

\begin{proof}
    Uniformly Cauchy in $\R$ $\implies$ uniformly continuous, follows from the general Stone-Weierstrass theorem.
\end{proof}

\begin{remark}
    Since $X$ is compact, then every continuous function is uniformly continuous. I.e. $\forall \varepsilon > 0, \ \exists \delta > 0$ s.t. 
    $$d_Y(f(x), f(x')) < \varepsilon, \quad \forall d(x, x') < \delta$$
\end{remark}


\begin{definition}
    $\mathcal{F} \subseteq C(X)$ is said to be equicontinuous iff $\forall \varepsilon > 0$, $\exists \delta > 0$ s.t. 
    $$d(f(x), f(x')) < \varepsilon, \quad \forall x, x' \ d(x,x') < \delta, \ \forall f \in \mathcal{F}$$
\end{definition}

\begin{theorem}[Arzela-Ascoli]
    $\mathcal{F} \subseteq C(X)$ is totally bounded $\iff$ $\mathcal{F}$ is bounded and equicontinuous.
\end{theorem}

\begin{proof}
    \begin{itemize}
        \item[$\implies$] $\mathcal{F}$ is bounded since it is totally bounded. Let $\varepsilon > 0$, since $\mathcal{F}$ is totally bounded, there exist $f_1, f_2, \cdots, f_N \in \mathcal{F}$ s.t. 
        $$\mathcal{F} \subseteq B(f_1, \varepsilon/3) \cup \cdots \cup B(f_N, \varepsilon/3)$$
        Now $f_i \in C(X)$, so uniformly continuous, and so $\exists \delta_i > 0$ s.t. 
        $$\abs{f_i(x) - f_i(y)} < \varepsilon/3, \quad \forall d(x,y) < \delta_i$$
        Let $\delta = \min_i \delta_i$. Now take $f \in \mathcal{F}$, and $x, y \in X$ with $d(x,y) < \delta$, now $\exists i_0 \in \set{1, \cdots, N}$ s.t. 
        $f \in B(f_{i_0}, \varepsilon/3)$. Hence 
        \begin{align*}
            \abs{f(x) - f(y)} &\le \abs{f(x) - f_{i_0}(x)} + \abs{f_{i_0}(x) - f_{i_0}(y)} + \abs{f_{i_0}(y) - f(y)} \\
            &\le 2 \times d_\infty(f, f_{i_0}) + \varepsilon/3 \\
            &< \varepsilon
        \end{align*} 
        Therefore $\mathcal{F}$ is equicontinuous.

        \item[$\impliedby$] Let $\varepsilon > 0$, then $\exists \delta > 0$ s.t $\forall f \in \mathcal{F}$ we have 
        $$\abs{f(x) - f(y)} < \varepsilon/3, \quad \forall x,y \in X \text{ s.t. } d(x,y) < \delta$$
        Now $X$ is compact, then we can cover $X$ with finitely many $\delta$-balls, i.e.
        $$X = B(x_1, \delta) \cup \cdots \cup B(x_N, \delta)$$
        Define a map $\varphi \colon \mathcal{F} \to \R^N$, 
        $$\varphi(f) = (f(x_1), \cdots, f(x_N))$$
        We want to show that $\varphi(\mathcal{F})$ is bounded in $\R^N$. We know that $\mathcal{F}$ is bounded in $C(X)$, so $\exists M > 0$ s.t. 
        $$\abs{f(x)} \le M, \quad \forall f \in \mathcal{F}, \ \forall x \in X$$
        Therefore 
        $$\abs{\varphi(f)} = \sqrt{f(x_1)^2 + \cdots + f(x_N)^2} \le \sqrt{N} M$$
        Therefore $\varphi(\mathcal{F})$ is bounded in $\R^N$. Therefore $\varphi(\mathcal{F})$ is totally bounded. Then 
        there exists $f_1, f_2, \cdots, f_K \in \mathcal{F}$ s.t. 
        $$\varphi(\mathcal{F}) \subseteq B(\varphi(f_1), \varepsilon/3) \cup \cdots \cup B(\varphi(f_K), \varepsilon/3)$$
        We want to show that 
        $$\mathcal{F} \subseteq B(f_1, \varepsilon) \cup \cdots \cup B(f_K, \varepsilon)$$
        Let $f \in \mathcal{F}$, then $\varphi(f) \in \varphi(\mathcal{F})$ i.e. $\exists i_{0}$ s.t. 
        $\varphi(f) \in B(\varphi(f_{i_0}), \varepsilon/3)$. Recall that 
        $$\varphi(f) = (f(x_1), \cdots, f(x_N))$$
        But notice that for each $j$, we have 
        $$\abs{f(x_j) - f_{i_0}(x_j)} < \varepsilon / 3, \quad \forall j =1, \cdots, N$$
        Let $x \in X$, then $x \in B(x_{j_0}, \delta)$ for some  
        \begin{align*}
            \abs{f(x) - f_{i_0}(x)} &\le \abs{f(x) - f(x_{j_0})} + \abs{f(x_{j_0}) - f_{i_0}(x_{j_0})} + \abs{f_{i_0}(x_{j_0}) + f_{i_0}(x)} \\
            &< \varepsilon
        \end{align*} 
        Since this holds for each $x \in X$ we have 
        $$\abs{f(x) - f_{i_0}(x)} < \varepsilon, \quad \forall x \in X$$
        Then since $X$ is compact
        $$\sup \abs{f(x) - f_{i_0}(x)} = \max \abs{f(x) - f_{i_0}(x)} < \varepsilon$$
        Hence $f \in B(f_{i_0}, \varepsilon)$, therefore $\mathcal{F}$ is totally bounded. 
    \end{itemize}
\end{proof}

\begin{corollary}
    $\mathcal{F} \subseteq C(X)$ is compact iff $\mathcal{F}$ is closed and bounded and equicontinuous.
\end{corollary}

\begin{lemma}
    $\mathcal{F} \subseteq C(X)$ is equicontinuous, then $\overline{\mathcal{F}}$ is equicontinuous.
\end{lemma}

\begin{proof}
    Take $\varepsilon > 0$ and $f \in \overline{f}$, then $\exists f_n \in \mathcal{F}$ s.t. 
    $$\sup_{X} \abs{f_n - f}$$
    but $\mathcal{F}$ is equicontinuous, so $\exists \delta > 0$ s.t. $\forall h \in \mathcal{F}$
    $$\abs{h(x) - h(x')} < \varepsilon/3, \quad \forall x, x' \text{ s.t. } d(x,x') < \delta$$
    for $d(x,x') < \delta$ 
    \begin{align*}
        \abs{f(x) - f(x')} &\le \abs{f(x) - f_n(x)} + \abs{f_n(x) - f_n(x')} + \abs{f_n(x') - f(x')} \\
        &\le 2 \sup_X \abs{f - f_n} + \varepsilon/3 
    \end{align*}
    But $\exists N$ s.t. for $n \ge N$ $\sup \abs{f - f_n} < \varepsilon/3$ and 
    $$\abs{f(x) - f(x')} < \varepsilon$$
\end{proof}

\begin{theorem}[Sequential Arzela-Ascoli]
    Let $f_n$ be a sequence in $C(X)$ s.t. 
    \begin{enumerate}[label = \roman*)]
        \item $f_n$ is uniformly bounded, i.e. $\abs{f_n(x)} \le M, \ \forall x \in X, \ \forall n \in \N$. 
        \item $f_n$ is equicontinuous, i.e. $\forall \varepsilon > 0, \ \exists \delta$ s.t. 
        $$\abs{f_n(x) - f_n(x')} < \varepsilon, \ \forall x, x' \ d(x,x') < \delta, \ \forall n \in \N$$
    \end{enumerate}
    Then $f_n$ has a subsequence that converges uniformly to a function $f \in C(X)$.
\end{theorem}

\begin{proof}
    $\mathcal{F} = \set{f_n \in \N}$, then $\overline{\mathcal{F}}$ is closed, bounded and equicontinuous. 
    Then $\overline{\mathcal{F}}$ is compact and hence is sequentially compact.
\end{proof}

\begin{example}
    In $C([0,1])$, let $\mathcal{F} = \set{f \in \operatorname{Lip}_1([0,1]) \colon f(0) = \alpha}$ where 
    $$\operatorname{Lip}_1([0,1]) = \set{f \colon \abs{f(x) - f(y)} \le \abs{x-y}, \ \forall x, y \in X}$$
    Now $\mathcal{F}$ is compact in $C(X)$. 
    \begin{itemize}
        \item \textbf{Closed.} Take $f_n \in \mathcal{F}$ s.t. $f_n \to f$ in $C([0,1])$. Clearly $f_n(0) = \alpha$ hence $f(0) = \alpha$. Now 
        $$\abs{f_n(x) - f_n(y)} \le \abs{x-y}$$
        let $n \to \infty$ we have 
        $$\abs{f(x) - f(y)} \le \abs{x-y}$$
        Hence $f \in \mathcal{F}$.

        \item \textbf{Bounded.} Now let $f \in \mathcal{F}$, then $\abs{f(x) - f(0)} \le \abs{x - \alpha} \le 1 + \alpha$, therefore for any $f \in \mathcal{F}$ and $x \in X$ we have
        $$\abs{f(x)} \le \abs{f(x) - f(0)} + \abs{f(0)} \le 1 + 2 \alpha$$
        Therefore 
        $$\sup \abs{f(x)} \le 1 + 2 \alpha$$
        So $\mathcal{F}$ is bounded.

        \item \textbf{Equicontinuous.} Let $\varepsilon > 0$, take $\delta = \varepsilon$, then for $\abs{x-y} < \delta = \varepsilon$ we have 
        $$\abs{f(x) - f(y)} \le \abs{x-y} < \varepsilon$$
    \end{itemize} 
    By Arzela-Ascoli, we get that $\mathcal{F}$ is compact.
\end{example}

\begin{example}
    Now fix $f \in C([0,1])$, and look at $m = \inf_{g} \set{d_\infty (f,g) \in \operatorname{Lip}_1 ([0,1]) \colon g(0) = f(0)}$, want to show that this is attained. 
    Indeed, let $g_n$ be a sequence in $\operatorname{Lip}_1([0,1])$ $g_n(0) = f(0)$ with $d_\infty(f, g_n) \to m$. By compactness, $g_n$ has convergent subsequence $g_{n_k}$ in $C(X)$ to a function 
    $g \in \operatorname{Lip}_1([0,1])$. Now $d_\infty(f,g) = \lim_{k \to \infty} d_\infty(f, g_{n_k}) = m$. 
\end{example}

\subsection{Baire Category Theorem}

\begin{definition}
    $(X, \tau)$ topological space, $A \subseteq X$. We say that $A$ is nowhere dense in $X$ iff $\overline{A}$ has empty interior.
\end{definition}

\begin{example}
    In $\R$, all singletons (or discrete sets) are nowhere dense in $\R$. But $\Q$ is not nowhere dense in $\R$.
\end{example}

\begin{proposition}
    $A$ is nowhere dense in $X$ $\iff X \setminus \overline{A}$ is dense. 
\end{proposition}

\begin{proof}
    \begin{itemize}
        \item[$\implies$] take $x \in X$ and $U_x \ni x$ open neighborhood, but $x$ is not an interior point of $\overline{A}$ because $\mathring{\overline{A}} = \varnothing$, therefore $U_x \cap \overline{A} \ne \varnothing$. 
        \item[$\impliedby$] Same. 
    \end{itemize}
\end{proof}


\begin{remark}
    $A$ closed set is nowhere dense iff $\mathring{A} = \varnothing$.   
\end{remark}

\begin{definition}
    $(X, \tau)$ topological space and $A \subseteq X$. We say that $A$ is first Category or meagre in $X$ iff $A$ can be written as a countable union of nowhere dense subsets of $X$.
    Otherwise we say that $A$ is second category (or non meagre).
\end{definition}

\begin{example}
    $\Q$ is first category in $\R$. 
\end{example}

\begin{remark}
    If a set $A$ is second Category with $A = \bigcup_{n=1}^\infty F_n$ with $F_n$ closed, then $\exists n_0$ with $F_{n_0}$ has nonempty interior. 
\end{remark}

\begin{theorem}[Baire Category Theorem]
    $(X, d)$ complete metric space, then if $G_1, G_2, \cdots$ are open dense subsets of $X$ then $\bigcap_{n=1}^\infty G_n$ is dense in $X$.
    Equivalently, every nonempty open subset of $X$ is second Category. In particular, if $X = \bigcup_{n=1}^\infty F_n$ with $F_n$ closed, then $\exists n_0$ with $F_{n_0}$ has nonempty interior.
\end{theorem}

\paragraph{Why equivalently?} 
\begin{itemize}
    \item[$\implies$] Let $G$ open in $X$, suppose that we can write $G = \bigcup_{n=1}^\infty A_n$ where $A_n$ is nowhere dense. Take $G_i = X \setminus \overline{A}$, 
    note that $G_i$ is open and dense in $X$, therefore $\bigcap_{n=1}^\infty G_n$ is dense in $X$. But 
    \begin{align*}
        \bigcap_{n=1}^\infty G_n &= \bigcap_{n=1}^\infty (\overline{A_n})^c \\
        &= \left(\bigcup_{n=1}^\infty \overline{A_n}\right)^c 
    \end{align*}
    Which is dense in $X$. But we have that 
    $$\bigcup_{n=1}^\infty \overline{A_n} \supseteq G$$
    with $G$ open, contradiction.

    \item[$\impliedby$] Let $G_1, G_2, \cdots$ be open dense subsets of $X$. Note that $G_i$ is second Category, and suppose $\bigcap_{n=1}^\infty G_n$ is not dense, 
    so $\exists H \subseteq \left(\bigcap_{n=1}^\infty G_n\right)^c = \bigcup_{n=1}^\infty G_n^c$ with $H$ open. But this means that 
    $$H = \bigcup_{n=1}^\infty G_n^c \cap H$$
    with $H$ second category, then $\exists n_0 \in \N$ and an open set $U \subseteq G_{n_0}^c \cap H \subseteq G_{n_0}^c$, contradiction.
\end{itemize} 


\begin{proof}[Proof of Baire Category Theorem]
    Take $x_0 \in X$ and $\delta > 0$, take $B_0 = B(x_0, \delta)$. Now $B_0$ intersects $G_1$, and $\exists x_1 \in B_0 \cap G_1$ and 
    $\delta_1 < \delta/2$ s.t. 
    $$B_1 = B(x_1, \delta_1) \subseteq \overline{B(x_1, \delta_1)} \subseteq B_0 \cap G_1$$
    Now $B_1$ intersects $G_2$, so $\exists x_2 \in B_1 \cap G_2$ and $\delta_2 < \delta/2^2$ with 
    $$B_2 = B(x_2, \delta_2) \subseteq \overline{B(x_2, \delta_2)} \subseteq B_1 \cap G_2 \subseteq B_0 \cap G_1 \cap G_2$$
    Continue this recursively to get $x_n \in B_{n-1} \cap G_n$ and $\delta_n < \delta/2^n$ s.t. 
    $$B_n = B(x_n, \delta_n) \subseteq \overline{B(x_n, \delta_n)} \subseteq B_0 \cap G_1 \cap \cdots \cap G_n$$
    We show that $x_n$ is a Cauchy sequence in $X$. Indeed, let $n,m > N$ for some $N \in \N$ then $x_n, x_m \in B_N = B(x_N, \delta_N)$, so 
    $$d(x_n, x_m) < 2 \delta/2^{N} = \delta/2^{N-1}$$
    Now let $\varepsilon > 0$, choose $N$ s.t. $\delta/2^{N-1} < \varepsilon$, then 
    $$d(x_n, x_m)< \varepsilon$$
    so $x_n$ is Cauchy. By completeness, $x_n$ converges to some $z \in X$. We show that $z \in \left(\bigcap_{n=1}^\infty G_n\right) \cap B_0$. 
    Fix $m$, we have  
    $$x_n \in \overline{B_m}, \quad \forall n > m$$ 
    Therefore $z \in \overline{B_m} \subseteq B_0 \cap \left(\bigcap_{k=1}^m G_k\right)$, since this holds for every $m$, we have 
    $$z \in B_0 \cap \left(\bigcap_{k=1}^\infty G_k\right)$$
\end{proof}

\section{Normed Spaces}

\begin{definition}
    Given a vector space $V$ over $\F$, a norm is a map 
    $$\norm{\cdot} \colon V \to \R$$
    satisfying 
    \begin{enumerate}[label = \roman*)]
        \item $\norm{x} \ge 0$ with equality iff $x = 0$.
        \item $\norm{\alpha x} = \abs{\alpha} \cdot \norm{x}$, $\alpha \in \F, \ x \in V$.
        \item $\norm{x + y} \le \norm{x} + \norm{y}, \ \forall x, y \in V$.
    \end{enumerate}
\end{definition}

\begin{example}
    $\ell^\infty(\R)$ (set of bounded sequences) under coordinate-wise addition and scalar multiplication. 
    $$\norm{x}_\infty = \sup_{i \in \N} \abs{x_i} = d_\infty (x, 0)$$
    $C(X)$ with $X$ compact 
    $$\norm{f}_\infty = \sup_{x \in X} \abs{f(x)}$$ 
\end{example}

\begin{proposition}
    Every norm induces a metric 
    $$d(x, y) = \norm{x-y}$$
    This allows us to deal with normed space as metric spaces w.r.t. induced metric.
\end{proposition}

\begin{definition}
    $V$ is a vector space and $\norm{\cdot}_a, \ \norm{\cdot}_b$ norms on $V$, we say that $\norm{\cdot}_a, \ \norm{\cdot}_b$ are equivalent iff 
    they induce the same topology.
\end{definition}

\begin{proposition}
    $V$ vector space, $\norm{\cdot}_a, \ \norm{\cdot}_b$ norms on $V$. The following are equivalent: 
    \begin{itemize}
        \item $\norm{\cdot}_a, \ \norm{\cdot}_b$ are equivalent.
        \item $\exists c_1, c_2 > 0$ s.t. 
        $$c_2 \norm{x}_b \le \norm{x}_a \le c_1 \norm{x}_b, \quad \forall x \in V$$
        \item $x_n \to x_0$ in $(V, \norm{\cdot}_a) \iff x_n \to x_0$ in $(V, \norm{\cdot}_b)$.
    \end{itemize}
\end{proposition}

\begin{proof}
    \begin{itemize}
        \item[$(i) \implies (ii)$] $B_{\norm{\cdot}_a} (0,1) = \set{x \in V, \ \norm{x}_a < 1}$ is open in $(V, \norm{\cdot}_a)$ then 
        $B_{\norm{\cdot}_a}(0,1)$ is open in $(V, \norm{\cdot}_b)$. Therefore $0$ is an interior point of $B_{\norm{\cdot}_a}(0,1)$ w.r.t. the $\norm{\cdot}_b$ norm.
        Therefore $\exists \delta > 0$ s.t. 
        $$B_{\norm{\cdot}_b} (0, \delta) \subseteq B_{\norm{\cdot}_a} (0,1)$$
        Now take $x \in V, \ x \ne 0$, then 
        $$\frac{\delta}{2} \frac{x}{\norm{x}_b} \in B_{\norm{\cdot}_b}(0, \delta) \subseteq B_{\norm{\cdot}_a} (0,1)$$
        Hence 
        $$\abs{\frac{\delta}{2} \frac{x}{\norm{x}_b}}_a < 1 \implies \norm{x}_b < \frac{2}{\delta} \abs{x}_b$$
        do the same by switching $\norm{\cdot}_a$ and $\norm{\cdot}_b$. 

        \item[$(ii) \implies (iii)$] We have that 
        $$c_1 \norm{x_n - x_0}_b \le \norm{x_n - x_0}_a \le c_1 \norm{x_n - x_0}$$

        \item[$(iii) \implies (i)$] Closed in $(V, \norm{\cdot}_a) \iff $ Closed in $(V, \norm{\cdot}_b)$ by sequential definition.
    \end{itemize}
\end{proof}

\begin{definition}
    Take $V$ finite dimensional vector space and $S = \set{v_1, \cdots, v_n}$ a basis of $V$. Take $x\in V$, then $\exists ! \ r_1, r_2, \cdots, r_n$ s.t. 
    $$x = r_1 v_1 + \cdots + r_n v_n$$
    We define 
    $$\norm{x}_\infty := \max_{1 \le i \le n} \abs{r_i}$$
\end{definition}

We then have the following: 
\begin{proposition}
    $V$ is a f.d. vector space and $S = \set{v_1, \cdots, v_n}$ a basis: 
    \begin{enumerate}
        \item $(V, \norm{\cdot}_\infty)$ is complete. 
        \item $\set{x \in V \colon \abs{x}_\infty \le M}$ is compact. 
        \item $K \subseteq (V, \norm{\cdot}_\infty)$ is compact iff $K$ is closed and bounded. 
    \end{enumerate}
\end{proposition}

\begin{proof}
    \begin{enumerate}
        \item Let $x_k$ be a Cauchy sequence in $V$, then 
        $$x_k = r_{k1} v_1 +\cdots + r_{kn} v_n$$
        Then 
        \begin{align*}
            \abs{r_{ki} - r_{li}} &\le \norm{x_k - x_l}_\infty 
        \end{align*}
        Then the sequence $(r_{ki})_k$ is Cauchy in $\R$ for every $i$, then $r_{ki}$ converges to some $r_i$. Take $x = r_1 v_1 + \cdots + r_n v_n$, we get that 
        \begin{align*}
            \norm{x_k - x}_\infty &= \norm{(r_{k1} - r-1) v_1 + \cdots + (r_{kn} - r_n) v_n}_\infty \\
            &\le \norm{r_{k1} - r_1} \norm{v_1}_\infty + \cdots + \norm{r_{kn} - r_n} \abs{v_n} \\
            &\to 0
        \end{align*}

        \item Take a sequence $x$ s.t. $\norm{x_k}_\infty \le M$. Apply Bolzano Weierstrass $n$ times and take subsequence of the subsequence.

        \item $(\implies)$ is always true. $(\impliedby)$ since $K$ is bounded, then $\exists M$ s.t. 
        $$K \subseteq \set{x \in V \colon \norm{x}_\infty \le M} =: Q$$
        But $K$ is closed $Q$ is compact, then $K$ is compact. 
    \end{enumerate}
\end{proof}

\begin{theorem}
    $V$ is a finite dimensional vector space, then all norms on $V$ are equivalent.
\end{theorem}

\begin{corollary}
    $V$ finite dimensional and $\norm{\cdot}$ is a norm on $V$ then $(V, \norm{\cdot})$ has the properties of the theorem above.

    Moreover, if $V$ is a normed space and $M$ is a finite dimensional subspace of $V$.
\end{corollary}

\begin{remark}
    Note that the converse is not true in general, take ${P}([a,b])$ as a subspace of $C([a,b])$.
\end{remark}

\begin{proof}
    Take $S = \set{v_1, \cdots, v_n}$ a basis of $V$, define 
    $$\norm{x}_\infty = \max_{1 \le i \le N} \abs{r_i}$$
    Let $\norm{\cdot}_a$ be a norm on $V$, we will show that $\norm{\cdot}_a$ is equivalent to $\norm{\cdot}$. 
    \begin{align*}
        \norm{x}_a &= \norm{r_1 v_1 + \cdots + r_n v_n}_a \\
        &\le \abs{r_1} \norm{v_1}_a + \cdots + \abs{r_n} \norm{v_n}_a \\
        &\le (\norm{v_1}_a + \cdots + \norm{v_n}_a) \norm{x}_\infty 
    \end{align*} 
    Let $\varphi \colon (V, \norm{\cdot}_\infty) \to \R$, note that this is continuous since if $x_n \to x_0$ in $\norm{\cdot}_\infty$, then 
    \begin{align*}
        \abs{\varphi(x_n) - \varphi(x_0)} &= (\norm{x_n}_a - \norm{x_0}) \\
        &\le \norm{x_n - x_0}_a \\
        &\le c_1 \norm{x_n - x_0}_\infty \\
        &\to 0
    \end{align*}
    Now let $K = \set{x \in V \colon \norm{x}_\infty = 1}$, note that $K$ is compact (closed and bounded), so 
    $\varphi(K)$ is compact in $\R$ so $m = \inf \varphi(K)$ is attained at some $x_0 \in K, \ x_0 \ne 0$. Then 
    $m = \norm{x_0}_a > 0$. So $\forall x \in V$ s.t. $\norm{x}_\infty = 1$ we have 
    $$\abs{x}_a \ge m$$
    Take $v \in V$ s.t. $x \ne 0$ we have that 
    \begin{align*}
        \norm{\frac{x}{\norm{x}\infty}}_a &\ge m \\
        \norm{x}_a \ge m \norm{x}_\infty
    \end{align*}
\end{proof}

\begin{lemma}[Riesz Lemma]
    $(V, \norm{\cdot})$ is a normed v.s. and $M$ a proper closed subspace of $V$. Then for every $0 < \alpha < 1$ there exists $x \in V$ s.t. 
    $\norm{x} = 1$ and 
    $$d(x, M) \ge \alpha$$
\end{lemma}

\begin{remark}
    If $x$ is s.t. $\norm{x} = 1$, then $d(x, M) \le \abs{x - 0} \le 1$.
\end{remark}

\begin{proof}
    Take $x_0 \not\in M$, $d(x_0, M) > 0$ since $M$ is closed. Fix $0 < \alpha < 1$, there exists $m_0 \in M$ s.t. 
    $$ d(x_0, M) \le \norm{x_0 - m_0} < \frac{d(x_0,M)}{\alpha}$$
    Define 
    $$x = \frac{x_0 - m_0}{\norm{x_0-m_0}}, \quad \norm{x} = 1$$
    Let $m \in M$, then 
    \begin{align*}
        \norm{x-m} &= \norm{\frac{x_0 - m_0}{\norm{x_0-m_0}} - m} \\
        &= \frac{\norm{x_0 - (m_0 + m \norm{x_0 - m_0})}}{\norm{x_0 - m_0}}
    \end{align*} 
    But $m_0 + m \norm{x_0 - m_0} \in M$, so 
    $$\norm{x-m} \ge \frac{d(x_0, M)}{\norm{x_0 - m_0}} > \alpha$$
    Therefore 
    $$d(x, M) = \inf_{m \in M} \norm{x-m} \ge \alpha$$
\end{proof}

\begin{theorem}
    $(V, \norm{\cdot})$ normed space, the closed unit sphere in $V$ is compact iff $V$ is finite dimensional.
\end{theorem}


\begin{proof}
    \begin{itemize}
        \item[$\impliedby$] ok. 
        
        \item[$\implies$] Assume that $V$ is infinite dimensional, then we can construct 
        $$E_1 = \Span \set{v_1}, \ E_2 = \Span \set{v_1, v_2}, \ \cdots$$
        With $E_i$ finite dimensional and $\set{v_i}$ linearly independent. Then 
        $$E_1 \subsetneq E_2 \subsetneq E_3 \subsetneq \cdots$$
        Fix some $u_1 \in E_1$ with $\norm{u_1} = 1$, now by Riesz' lemma $\exists u_2 \in E_2$ s.t. $\norm{u_2} = 1$ and 
        $d(u_2, E_1) \ge 1/2$. Now $\exists u_3 \in E_3$ with $\norm{u_3} = 1$ and $d(u_3, E_2) \ge 1/2$, recursively we can construct $u_n \in E_n$ s.t. 
        $\norm{u_n} = 1$ and $d(u_n, E_{n-1}) \ge 1$. Now notice that for $m \ne n$ we have 
        \begin{align*}
            \norm{u_m - u_n} \ge 1/2
        \end{align*}
        Therefore $u_n$ is a sequence which belongs to the unit sphere that has no Cauchy (and hence convergent) subsequence. So the unit sphere is not sequentially compact, and hence not compact.
    \end{itemize}
\end{proof}

\begin{definition}
    A complete normed space is called a Banach space.
\end{definition}

\begin{example}
    Some Banach spaces: $\R^N$.
    $\ell^\infty(\R)$ (which is the set of bounded sequences). 
    $\ell^(\R) = \set{x \colon (x_1, x_2, \cdots,) \colon \sum \abs{x_i} < \infty}$. 
    $C(X)$ when $X$ is compact.

    $P(\R)$ is not Banach.
\end{example}

\textbf{$\ell^p$ Spaces.} 

\begin{lemma}
    $1 < p,q < \infty$ s.t. $\frac{1}{p}+ \frac{1}{q} = 1$ and $x,y > 0$ s.t. $xy = 1$ then 
    $$\frac{x^p}{p} + \frac{y^q}{q} \ge 1$$
\end{lemma}

\begin{proof}
    $f(x,y) = \frac{x^p}{p} + \frac{y^q}{q}$, we want to minimize $f(x,y)$ on $g(x,y) = xy = 1$. Note that this minimum exists since this is a closed set and $f(x,y) \ge 0$ 
    for all $x, y$ s.t. $xy = 1$. Using Lagrange multipliers, we solve 
    $$
    \begin{cases}
        \nabla f &= \lambda \nabla g \\
        xy &= 1
        \end{cases}
    $$
    i.e.
    $$\begin{cases}
            x^{p-1} = \lambda y \\
            y^{q-1} = \lambda y \\
            xy = 1
        \end{cases}
    $$
    Then we get 
    $$x^p = \lambda xy = \lambda, \quad y^q = \lambda$$
    Then $x = \lambda^{1/p}, y = \lambda^{1/q}$, then 
    $$\lambda^{\frac{1}{p} + 1/q} = 1 \implies \lambda = 1$$
    Then $x = y = 1$. So the minimum is $f(1,1 ) = 1$.
\end{proof}

\begin{corollary}
    $\frac{1}{p} + \frac{1}{q} = 1$, $x,y \ge 0$ then 
    $$\frac{x^p}{p} + \frac{y^q}{q} \ge xy$$
\end{corollary}

\begin{proof}
    Take $X = \frac{x}{(xy)^{1/p}}$ and $Y = \frac{y}{(xy)^{1/q}}$ then 
    $$XY = \frac{xy}{(xy)^{1/p + 1/q}} = 1$$
    Then apply the lemma, we get 
    \begin{align*}
        \frac{X^p}{p} + \frac{Y^q}{q} &\ge 1 \\
        \frac{x^p}{(xy)p} + \frac{y^q}{(xy)q} &\ge 1\\
        \frac{x^p}{p} + \frac{y^q}{q} &\ge xy  
    \end{align*}
\end{proof}

\begin{theorem}[Young's inequality]
    Let $x,y \in \R$ then 
    $$\abs{xy} \le \frac{\abs{x}^p}{p} + \frac{\abs{y}^q}{q}$$
    with $\frac{1}{p} + \frac{1}{q} = 1$. 
\end{theorem}

\begin{theorem}[Holder inequality]
    Let $x_1, x_2, \cdots, x_n \in \R$ and $y_1, y_2, \cdots, y_n \in \R$ then 
    $$\sum_{i=1}^n \abs{x_iy_i} \le \left(\sum_{i=1}^n \abs{x_i}^p\right)^{1/p} \left(\sum_{i=1}^n \abs{y_i}^{q}\right)^{1/q}$$
    With $\frac{1}{p} + \frac{1}{q} = 1$.

    Letting $n \to \infty$ we have 
    $$\sum_{n=1}^\infty \abs{x_n y_n} \le \left(\sum_{n=1}^\infty \abs{x_n}^p\right)^{1/p} \left(\sum_{n=1}^\infty \abs{y_n}^q\right)^{1/q}$$
\end{theorem}

\begin{proof}
    Wlog suppose not all zeroes. Let $X_i = \frac{\abs{x_i}}{\left(\sum_{k=1}^n \abs{x_k}^p\right)^{1/p}}$ and $y_i = \frac{\abs{y_i}}{\left(\sum_{k=1}^n \abs{y_k}^q\right)^{1/q}}$. 
    From Young's inequality, we have that 
    \begin{align*}
        X_i Y_i &= \frac{\abs{x_i} \abs{y_i}}{\left(\sum_{k=1}^n \abs{x_k}^p\right)^{1/p} \left(\sum_{k=1}^n \abs{y_k}^q\right)^{1/q}} \\
        &\le \frac{X_i^p}{p}  + \frac{Y_i^q}{q} \\
        &= \frac{\abs{x_i}}{\left(\sum_{k=1}^n \abs{x_k}^p\right)p} + \frac{\abs{y_i}^q}{\left(\sum_{k=1}^n \abs{y_k}^q\right)q}
    \end{align*}
    Therefore 
    \begin{align*}
        \sum_{i=1}^n X_i Y_i &\le \frac{1}{p} + \frac{1}{q}\\
        &= 1
    \end{align*}
    Substituting we get the desired inequality.
\end{proof}

\begin{theorem}[Minkowski inequality]
    $\left(\abs{x_i + y_i}\right)^p \le \left(\sum_{i=1}^n \abs{x_i}^p\right)^{1/p} + \left(\sum_{i=1}^n \abs{y_i}^p\right)^p$. In the infinite case we get 
    $$\left(\sum_{n=1}^\infty \abs{x_n + y_n}^p\right)^{1/p} \le \left(\sum_{n=1}^\infty \abs{x_n}^p\right)^{1/p} + \left(\sum_{n=1}^\infty \abs{y_n}^p\right)^{1/p}$$
\end{theorem}

\begin{proof}
    If $p =1$, trivial. Take $p > 1$ and let $q$ be s.t. $\frac{1}{p} + \frac{1}{q} = 1$. Then 
    \begin{align*}
        \sum_{i=1}^n \abs{x_i + y_i}^p &= \sum_{i=1}^n \abs{x_i + y_i} \cdot \abs{x_i + y_i}^{p-1} \\
        &\le \sum_{i=1}^n \left(\abs{x_i} + \abs{y_i}\right) \abs{x_i + y_i}^{p-1} \\
        &= \sum_{i=1}^n \abs{x_i} \abs{x_i + y_i}^{p-1} + \sum_{i=1}^n \abs{y_i} \abs{x_i + y_i}^{p-1} \\
        &\le \left(\sum_{i=1}^n \abs{x_i}\right)^{1/p} \left(\sum_{i=1}^n \abs{x_i + y_i}^{(p-1)q}\right)^{1/q} \\
        &\quad + \left(\sum_{i=1}^n \abs{y_i}^p\right)^{1/p} \left(\sum_{i=1}^n \abs{x_i + y_i}^{(p-1)q}\right)^{1/q} 
    \end{align*}
    We have that $\frac{1}{p} + \frac{1}{q} = 1 \implies p+q = pq \implies p = q(p-1)$, therefore 
    \begin{align*}
        \sum_{i=1}^n \abs{x_i + y_i}^p &\le \left(\sum_{i=1}^n \abs{x_i + y_i}^p\right)^{1/q} \left(\left(\sum_{i=1}^n \abs{x_i}^[]\right)^{1/p} + \left(\sum_{i=1}^n \abs{y_i}^p\right)^{1/p}\right) \\
        (\sum_{i=1}^n \abs{x_i + y_i}^p)^{1-1/q = 1/p} &\le \left( \sum_{i=1}^n \abs{x_i}^p\right)^{1/p} + \left(\sum_{i=1}^n \abs{y_i}^p\right)^{1/p}
    \end{align*}
    as desired.
\end{proof}

\begin{definition}
    Let $p \ge 1$, we define the space $\ell^p(\R) = \set{x = (x_1 ,\cdots, ) \colon \sum_{n=1}^\infty \abs{x_n}^p < \infty}$ with the following norm 
    $\norm{x}_{\ell^p} = \left(\sum_{n=1}^\infty \abs{x_n}^p\right)^{1/p}$. 
\end{definition}

\begin{proposition}
    \begin{itemize}
        \item $\ell^p(\R)$ is complete $\forall 1 \le p \le \infty$. 
        \item $\ell^p(\R)$ is separable $\forall 1 \le p < \infty$ (recall that $\ell^\infty(\R)$ is not separable). Since $c_{00}(\Q)$ is dense in $\ell^p(\R)$.
        \item $K \subseteq \ell^p(\R)$ ($p \le 1 < \infty$) is compact iff it is closed and bounded and 
        $$\lim_{N \to \infty} \sup_{x \in K} \sum_{N=1}^\infty \abs{x_n}^p = 0$$
        \item $1 \le r \le s < \infty$, then 
        $$\ell^r(\R) \subseteq \ell^s(\R)$$
        Hence 
        $$\ell^1(\R) \subseteq \ell^r (\R) \subseteq \ell^s(\R) \subseteq \ell^\infty (\R)$$
        \item If $1 \le r < s < \infty$ then $\ell^r (\R) \subsetneq \ell^s(\R)$ 
        $$x_n = \frac{1}{n^{1/r}}$$
        Then $\sum \abs{x_n} < \infty$ but $\sum \abs{x_n}^r = \infty$ so 
        $x_n \in \ell^s(\R)$ and $x_n \not\in \ell^r(\R)$. 
    \end{itemize}
\end{proposition}

\subsection{Bounded Linear Transformation}

\begin{definition}
    $V,W$ vector spaces and $\mathcal{L}(V,W)$ is the space of linear transformations $T \colon V \to W$, this is a vector space.
\end{definition}

\begin{remark}
    If $V,W$ are finite dimensional then $\mathcal{L}(V,W)$ is finite dimensional with 
    $$\operatorname{dim} \mathcal{L}(V,W) = \operatorname{dim}(V) \cdot \operatorname{dim} (W)$$
\end{remark}

\begin{definition}
    Take $(V, \norm{\cdot}_V)$ and $(W, \norm{\cdot}_W)$ and $T \in \mathcal{L}(V,W)$ then we say that $T$ is bounded iff $\exists M > 0$ s.t. 
    $$\norm{Tx}_W \le M \norm{x}_V$$
    Equivalently $\exists M$ s.t. 
    $$\norm{Tx} \le M, \quad \forall x \text{ s.t. } \norm{x}_V = 1$$
\end{definition}

\begin{proposition}
    $(V, \norm{\cdot}_V)$ and $(W, \norm{\cdot}_W)$ normed spaces and $T \in \mathcal{L}(V,W)$, then 
    $$ T \text{ bounded} \iff  T \text{ is continuous} \iff T \text{ is continuous at } 0$$
\end{proposition}

\begin{proof}
    \begin{itemize}
        \item[$(i) \implies (ii)$] Then $\exists M$ s.t. $\norm{Tx} \le M \norm{x}$. Then 
        \begin{align*}
            \norm{Tx - Ty} &= \norm{T(x-y)} \\
            &\le M \norm{x-y}
        \end{align*}
        Then $T$ is Lipschitz so it is continuous. 

        \item[$(ii) \implies (iii)$] ok.
        
        \item[$(iii) \implies (i)$] Let $T$ be continuous at $0$, then $\exists \delta > 0$ s.t. $\norm{Tx - T0} < 1$ for $\norm{x - 0} < \delta$ but this gives 
        $$\norm{Tx} \le 1$$
        for $\norm{x} < \delta$. Now take $x \ne 0$, then 
        $$\norm{T \left(\frac{\delta}{2}\frac{x}{\norm{x}}\right)} \le 1$$
        So we get 
        $$\norm{Tx} \le \frac{\delta}{2} \norm{x}$$
    \end{itemize}
\end{proof}

\begin{remark}
    If $V$ is finite dimensional, and $T \in \mathcal{L}(V,W)$, then $T$ is continuous. Take $x \in V$, and $S = \set{v_1, \cdots, v_n}$ basis of $V$. Then 
    $$x = r_1 v_1 + \cdots + r_n v_n \implies \norm{Tx} \le r_1 \norm{Tv_1} + \cdots + r_n \norm{Tv_n} \le (\norm{Tv_1} + \cdots \norm{Tv_n}) \norm{x}_\infty$$
    But all norms on $V$ are equivalent, then 
    $$\norm{Tx}_W \le M \norm{x}_V$$
    Then $T$ is bounded.  
\end{remark}

\begin{definition}
    $(V, \norm{\cdot}), \ (W, \norm{\cdot})$ normed spaces, we define $B(V,W)$ to be the space of bounded linear transformations with the following norm
    $$\norm{T}_{B(V,W)} = \sup_{\substack{x \in V \\ \norm{x} = 1}} \norm{Tx}_W = \sup_{\substack{x \in V \\ x \ne 0}} \frac{\norm{Tx}}{\norm{x}}$$
    Check that this is subspace of $\ell(V, W)$ and $\norm{\cdot}_{B(V,W)}$ is a norm.
\end{definition}

\begin{remark}
    Notice that for $x \in V$ $x \ne 0$ we have that 
    \begin{align*}
        \abs{T \frac{x}{\norm{x}}} &\le \norm{T} \\
        \norm{Tx} \le \norm{T} \cdot \norm{x}
    \end{align*}
\end{remark}

\begin{proposition}
    If $W$ is Banach, then $B(V, W)$ is Banach.
\end{proposition}

\begin{proof}
    Let $T_n$ be a Cauchy sequence in $B(V,W)$, note that 
    \begin{align*}
        \norm{T_n x - T_m x} &\le \norm{T_n - T_m} \cdot \norm{x}
    \end{align*}
    But $T_m$ is Cauchy in $B(V,W)$, therefore $T_n(x)$ is Cauchy is $W$ (which is Banach), then in converges. So we define 
    $$T \colon V \to W, \quad T(x) = \lim_{n \to \infty} T_n(x)$$

    We want to show $T$ is linear, this is trivial by properties of limits and that $T_n$ are linear. Moreover $\set{T_n}$ is Cauchy 
    in $B(V,W)$, then it is bounded i.e. $\exists M > 0$ s.t. $\norm{T_n} \le M, \ \forall n$. Now 
    \begin{align*}
        \norm{T_nx} &\le \norm{T_n} \cdot \norm{x} \\
        \norm{T_nx} &\le M \norm{x}
    \end{align*} 
    Letting $n \to \infty$ we get that $\norm{Tx} \le M \norm{x}$ so $T \in B(V,W)$.
    
    We want to show that $T_n \to T$ in $B(V,W)$. Given $\varepsilon > 0$, since $T_n$ is Cauchy, so $\exists N_1$ s.t. 
    $$\norm{T_n - T_m} < \varepsilon/2, \quad \forall n,m \ge N_1$$
    Now fix $x \in V$, since $T_n x \to T x$ in $W$ we have that $\exists N_2(x) > N_1$ s.t. 
    $$\norm{T_{N_2} x - Tx} < \varepsilon/2$$
    Now for $n \ge N_1$ and $\norm{x} = 1$ we have  
    \begin{align*}
        \norm{T_nx - Tx} &< \norm{T_n x - T_{N_2} x} + \norm{T_{N_2} x - Tx} \\
        &\le \norm{T_n - T_{N_2}} \cdot \norm{x} + \frac{\varepsilon}{2} \\
        &< \varepsilon
    \end{align*}
    Therefore given $\varepsilon > 0$, $\exists N_1$ s..t $\forall x \in V$ with $\norm{x} = 1$ then 
    $$\norm{T_n x - Tx} < \varepsilon, \quad \forall n \ge N_1$$
    Therefore 
    $$\sup_{\substack{x \in V \\ \norm{x} = 1}} \norm{T_nx - T x} \le \varepsilon, \quad \forall n \ge N_1$$
    Hence 
    $$T_n \to T \in B(V, W)$$
\end{proof}

\begin{example}
    Define 
    $$T \colon \left(C([0,1]), \norm{\cdot}_\infty\right) \to \R, \quad T (f) = \int_{0}^1 f(x) \ dx$$
    Show that $T \in B(C([0,1]), \R)$ and find $\norm{T}$.

    \textbf{Solution:} Clearly $T$ is linear because integration is linear. Now 
    \begin{align*}
        \abs{Tf} &= \abs{\int_0^1 f(x) \ dx} \\
        &\le \int_0^1 \abs{f(x)} \ dx \\
        &\le \norm{f}_\infty
    \end{align*}
    Take $T \in B(C([0,1]), \R)$, note that 
    $$\sup_{\norm{f}_\infty} \norm{Tf} \le 1$$
    moreover if we take $f = 1$ we have that $\norm{Tf} = 1$. Therefore 
    $$\norm{T} = 1$$
\end{example}

\subsection{Duality} 

\begin{definition}
    Let $V$ normed space over $\F \in \set{\R, \C}$. Then the dual set $V^* = B(V,\F)$, which is the set of all bounded linear maps from $V \to \F$.
    In particular, the dual norm 
    $$\norm{f}_* = \sup_{\substack{x \in V \\ \norm{x} = 1}} \norm{f(x)} = \sup_{\substack{x \in V \\ x \ne 0}} \norm{\frac{\abs{f(x)}}{\norm{x}}}$$
\end{definition}

\begin{remark}
    If $V$ is finite dimensional, then $V^* = B(V, \F) = \mathcal{L}(V,\F)$, then 
    $$\dim V^* = \dim V < \infty$$
    Therefore $V^*$ and $V$ are isomorphic. 
\end{remark}

\textbf{notation.} Given $(V, \norm{\cdot}), \ (W, \norm{\cdot})$ we write $V \cong W$ iff $V$ and $W$ are isometrically isomorphic. i.e. $\exists \varphi \colon V \to W$ and isomorphism and isometry.

\begin{example}
    $(\ell^1 (\R))^* \cong \ell^\infty(\R)$, more particularly if $1 < p < q < \infty$ and $1/p + 1/q = 1$, then 
    $$(\ell^p)^* \cong \ell^q$$
    So for $p = 2$ we get that 
    $$(\ell^2)^* \cong \ell^2$$
    Now if $1 < p < q < + \infty$  then 
    $$(\ell^p)^{**} \cong (\ell^q)^* \cong (\ell^p)^*$$
    If $p = 1$, then 
    $$(\ell^1)^* \cong \ell^\infty \text{ but } (\ell^\infty)^* \ne \ell^1$$
\end{example}

\begin{theorem}
    $(\ell^1)^* \cong \ell^\infty$, moreover if $1 < p < q < \infty$ then 
    $$(\ell^p)^* \cong \ell^q$$
    Where $\frac{1}{p} + \frac{1}{q} = 1$
\end{theorem}

\begin{lemma}
    Let $1 \le p < \infty$ and let $x = (x_1, x_2, \cdots) \in \ell^p(\R)$ then 
    $$x = \sum_{n=1}^\infty x_n e_n = \lim_{N \to \infty} \sum_{n=1}^N x_n e_n$$
\end{lemma}

\begin{proof}
    Indeed note that 
    $$\sum_{n=1}^N x_n e_n = (x_1, x_2, \cdots, x_N, 0, 0, \cdots)$$
    and since $x_n \in \ell^p(\R)$ we have that 
    $$\sum_{n=1}^\infty \abs{x_n}^p < \infty$$
    Now 
    \begin{align*}
        \norm{x - \sum_{n=1}^N x_n e_n}_{\ell^p} &= \norm{(0,\cdots, 0, x_{N+1}, x_{N+2})}_p \\
        &= \left(\sum_{n=N+1}^\infty \abs{x_n}^p \right)^{1/p} \\
        &\to 0
    \end{align*}
\end{proof}

\begin{remark}
    Now for $\ell^\infty$ the result is false. Indeed consider 
    $$\1 = (1, 1, 1, \cdots)$$
    Then 
    $$\sum_{n=1}^N x_n e_n = (1,1,\cdots, 1, 0,0,\cdots)$$
    But 
    \begin{align*}
        \norm{x- \sum_{n=1}^N x_n e_n}_\infty &= \norm{(0,\cdots, 0, 1,1,\cdots)}_\infty \\
        &= 1
    \end{align*}
    Which means that 
    $$\lim_{N \to \infty} \sum_{n=1}^N x_n e_n \ne x$$
\end{remark}

\begin{proof}[Proof of $(\ell^1)^* \cong \ell^\infty$]
    Let $T \in (\ell^1)^*$, then $T \colon \ell^1 \to \R$ bounded. Let $x \in \ell^1$, then by the lemma we have that $x = \sum_{n=1}^\infty x_n e_n$. With $T$ continuous and linear therefore 
    $$Tx = \sum_{n=1}^\infty x_n Te_n$$
    Let $b = (Te_1, Te_2, \cdots)$, we want to show that $b \in \ell^\infty$. Indeed, since $T$ is bounded we have 
    $$\norm{Te_i} \le \norm{T} \cdot \norm{e_i} = \norm{T}$$
    Therefore 
    $$\sup_{i \in \N} \abs{Te_i} \le \norm{T}$$
    Hence $b \in \ell^\infty$ with 
    $$\norm{b}_\infty \le \norm{T}$$
    On the other hand 
    \begin{align*}
        \abs{Tx} &= \abs{x_n Te_n} \\
        &\le \sum_{n=1}^\infty \abs{x_n} \cdot \abs{b_n} \\ 
        \norm{Tx} &\le \norm{b}_\infty \norm{x}
    \end{align*}
    Therefore 
    $$\norm{T} \le \norm{b}_\infty$$
    Therefore 
    $$\norm{T}_* = \norm{b}_\infty$$
    Using all of this we know that the linear map 
    $$\psi \colon (\ell^1)^* \to \ell^\infty, \quad \psi(T) = (Te_1, Te_2, \cdots)$$
    is an isometry. To prove surjectivity, we can take $b = (b_1, b_2, \cdots) \in \ell^\infty$ and define $T_b \colon \ell^1 \to \R$ by 
    $$T_b x = \sum_{n=1}^\infty x_n b_n$$
    Note that this converges since 
    $$\sum_{n=1}^\infty \abs{x_n b_n} \le \norm{b} \norm{x}$$
    Then $T_b$ is linear and
    $$\abs{T_bx} \le \norm{b}_\infty \norm{x} \implies T_b \in (\ell^1)^*$$
    Then 
    \begin{align*}
        \psi(T_b) &= (T_be_1, T_b e_2, \cdots) \\
        &= (b_1, b_2, \cdots) \\
        &= b
    \end{align*}
\end{proof}

\section{Hahn-Banach Theorems}

\textbf{Motivation.} Let $(V, \norm{\cdot})$ be a normed space over $\F$ and let $\mathcal{M}$ be a subspace of $V$ and $T \colon M \to \F$ a bounded linear map. 
Notice that 
$$\norm{Tx - Ty} \le \norm{T} \cdot \norm{x - y}$$
Therefore $T$ is uniformly continuous and $\F$ is complete, then by the homework the map 
$$\tilde{T} \colon \overline{\mathcal{M}} \to \F, \quad \tilde{T} x = \lim_{n \to \infty} Tx_n, \ x_n \text{ a sequence in $\mathcal{M}$ converging to $x$}$$
is a continuous extension of $T$ into $\overline{\mathcal{M}}$. Moreover, $\mathcal{M}$ is a subspace then $\overline{\mathcal{M}}$ is also a subspace. Then $\tilde{T}$ is also 
linear. Then $\tilde{T} \in B(\overline{\mathcal{M}}, \F)$. Now by definition 
\begin{align*}
    \norm{\tilde{T}} &= \sup_{\substack{x \in \overline{\mathcal{M}} \\ x \ne 0}}\frac{\norm{\tilde{T} x}}{\norm{x}} \\
    &\ge \norm{T}
\end{align*}
Let $x \in \overline{\mathcal{M}}$ and $x_n \to x$ with $x_n \in \mathcal{M}$. Then 
\begin{align*}
    \abs{Tx_n} &\le \norm{T} \cdot \norm{x_n} 
\end{align*} 
Letting $n \to \infty$ we get 
$$\norm{\tilde{T} x} \le \norm{T} \cdot \norm{x} \implies \norm{\tilde{T}} \le \norm{T}$$

\begin{lemma}
    $(V, \norm{\cdot})$ normed space and $\mathcal{M}$ proper subspace of $V$ and $x_0 \not\in \mathcal{M}$. Let 
    $$f \colon \mathcal{M} \to \R$$ 
    a bounded linear map and let $x_0 \not\in \mathcal{M}$. Then there exists an extension 
    $$F \colon \mathcal{M} \oplus \Span \set{x_0} \to \R$$
    s.t. $F$ is a bounded linear map and 
    $$F_{|M} = f$$
    and 
    $$\norm{F} = \norm{f}$$
\end{lemma}

\begin{proof}
    Take $x \in \mathcal{M} \oplus \Span \set{x_0}$, then we can write 
    $$x = m + tx_0, \quad t \in \R$$
    We would like to have 
    \begin{align*}
        F(x) &= F(x + tx_0) \\
        &= f(m) + t c
    \end{align*}
    So we need to define $F(x_0) := c \in \R$ s.t. 
    $$\abs{F(m + tx_0)} \le \norm{f} \norm{m + t x_0} (*)$$
    Since this would give us that 
    $$\norm{F} = \norm{f}$$
    Clearly we have that $(*)$ holds for $t = 0$. 
    Note that it is enough to prove this for $t=1$, i.e. that 
    $$\abs{F(m + x_0)} \le \norm{f} \norm{m + x_0}$$
    Indeed, suppose that this is the case, then for $t \ne 0$
    \begin{align*}
        \abs{F(m + tx_0)} &=  \abs{t} \abs{F(\frac{m}{t} + x_0)} \\
        &\le \abs{t} \norm{f} \norm{\frac{m}{t} + x_0} \\
        &= \norm{f} \norm{m + tx_0}
    \end{align*} 
    Then it is enough to find $c$ s.t. 
    $$\abs{F(m + x_0)} \le \norm{f} \norm{m + x_0}$$
    Which is equivalent to saying that 
    $$-f(m) - \norm{f} \norm{m + x_0} \le c \le \norm{f} \norm{m + x_0} - f(m), \quad \forall m \in \mathcal{M}, \ (**)$$
    Take $m, m' \in \mathcal{M}$ then 
    \begin{align*}
        f(m') - f(m) &= f(m' - m) \\
        &\le \abs{f(m' - m)} \\
        &\le \norm{f} \norm{m' - m} \\
        &\le \norm{f} \left(\norm{m' + x_0} + \norm{m + x_0}\right)
    \end{align*}
    Which is equivalent to saying that 
    \begin{align*}
        f(m) - \norm{f} \norm{m + x_0} &\le - f(m') + \norm{f} \norm{m' + x_0}, \quad \forall m,m' \in \mathcal{M}
    \end{align*}
    Hence 
    $$\sup_{m} \left(-f(m) - \norm{f} \norm{m+x_0}\right) \le \inf_{m} \left(-f(m) + \norm{f} \norm{m+x_0}\right)$$
    And so we can find $c$ satisfying $(**)$. Then we can define 
    $$F(m + tx_0) = f(m) + tc_0$$
    Hence 
    \begin{align*}
        \abs{F(m + tx_0)} &\le \norm{f} \norm{m + tx_0} \\
    \end{align*}
    Therefore $F$ is bounded and clearly $\norm{F} = \norm{f}$ (this gives us $\le$ but from the definition $\ge$ is clear). 
\end{proof}

\subsection{Partially Ordered Sets} 

\begin{definition}
    Given a set $X$, a partial order on $X$ is a relation $\preceq$ on $X$ that is 
    \begin{enumerate}[label = \roman*)]
        \item Reflexive: $x \preceq x$. 
        \item Antisymmetric: $x \preceq y \text{ and } y \preceq x \implies x = y$. 
        \item Transitive: $x \preceq y \text{ and } y \preceq z \implies x \preceq z$.
    \end{enumerate}
    Then $(X, \preceq)$ is a partially ordered set.
\end{definition}

\begin{example}
    We can say that $(\mathcal{P}(X), \subseteq)$ is a partially ordered set.
    
    We can also define on $\N$ the relation 
    $$a \le b \iff a \text{ divides b}$$
    Notice that in this case $2 \le 6$ but $2$ and $3$ are not comparable.
\end{example}

\begin{definition}
    Given a partially ordered set $(X, \le)$ we say that $(X, \le)$ is totally ordered iff $\forall x, y \in X$ we have that 
    $x \le y \text{ or } y \le x$.
\end{definition}

\begin{definition}
    Given a partially ordered set $(X, \le)$, we say that $M \in X$ is a maximal element of $X$ iff 
    $\forall a \in X$ s.t. $a \ge M$ then $a = M$. 

    Let $A \subseteq X$ we say that $M \in X$ is an upper bound of $A$ iff $\forall a \in A$ we have 
    $a \le M$.  
\end{definition}

\begin{example}
    $\set{2,3,5, 30}$ with divisibility, then $30$ is an upper bound for this set. But if we take $A  =\set{2,3,5,6}$ as a subset of $\N$, then $30$ 
    is an upper bound of $A$, but the maximal elements of $A$ are $5,6$. 
\end{example}

\begin{definition}
    Let $(X, \le)$, be a partially ordered set. Then a chain $C$ of $X$ is totally ordered subset of $X$.
\end{definition}

\begin{theorem}[Zorn's Lemma]
    Let $(X, \tau)$ be a nonempty partially ordered set such that  every chain in $X$ has an upper bound. Then $X$ has a maximal element. 
\end{theorem}

\begin{theorem}
    Every vector space $V$ has a (Hamel) basis.
\end{theorem}

\begin{proof}
    The empty set is a basis of the zero vector space. 
    Suppose that $V \ne 0$, let $X$ be set the set of linearly independent subsets of $V$ ordered by inclusion. First notice that $X$ is nonempty 
    because for any $v_0 \ne 0$ in $V$ we have that $\set{v_0}$ is linearly independent so $\set{v_0} \in X$. 
    Let $\mathcal{C} = \set{C_\alpha}_{\alpha \in I}$ be a chain in $X$, then take 
    $$C = \bigcup_{\alpha \in I} C_\alpha$$
    We want to show that $C \in X$, i.e. that $C$ is L.I. Take $v_1, v_2, \cdots, v_n \in C$ with 
    $$r_1 v_1 + r_2 v_2 + \cdots + r_n v_n = 0$$
    But $v_i \in C_{\alpha_i}$ for some $\alpha_i \in I$, since $C_{\alpha_i}$ are L.I, then by renumbering we can say 
    $$C_{\alpha_1} \le C_{\alpha_2} \le \cdots \le C_{\alpha_n}$$
    Therefore $v_1, v_2, \cdots, v_n \in C_{\alpha_n}$ hence 
    $$r_1 = r_2 = \cdots = r_n = 0$$
    Hence $C$ is an upper bound of $\mathcal{C}$. 

    Then by Zorn's lemma $X$ has a maximal element $S$, so $S$ is linearly independent, we still need to show $\Span S = V$. Suppose it is not, $\exists w \in V \setminus \Span S$, 
    then $S \cup \set{w}$ is linearly independent and $S < S \cup \set{w}$, contradiction. Therefore $S$ is a basis of $V$.  
\end{proof}

\begin{theorem}[Hahn-Banch]
    Let $V$ normed vector space over $\R$ and $M$ supsace of $V$, and $f \colon M \to \R$ is a bounded linear functional. Then $\exists F \colon V \to \R$ bounded linear functional 
    such that $F_{|M} = f$ and $\norm{F} = \norm{f}$
\end{theorem}

\begin{proof}
    Let $X = \set{(h,D(h))}$ Where $D(h)$ is a subspace of $V$ containing $M$ and $h \colon D(h) \to \R$ a bounded linear functional, $h_{|M} = f$ and $\norm{h} = \norm{f}$ with the following partial ordered 
    $$(h_1, D(h_1)) \le (h_2, D(h_2)) \iff D(h_1) \subseteq D(h_2) \text{ and } {h_2}_{D(h_1)} = h_1$$
    Now clearly $X \ne \varnothing$ since $(f, M) \in X$. Now let $\mathcal{C} = \set{(h_\alpha, D(h_\alpha))}$ be a chain in $X$. Define 
    $$D = \bigcup_{\alpha \in I} D(h_\alpha)$$
    Then $M \subseteq D$, $D$ is a subspace of $V$ since for$ x,y \in D$ and $\lambda \in \R$ we have $x \in D(h_{\alpha_x})$ and $y \in D(h_{\alpha_y})$ but $\mathcal{C}$ is a chain so 
    wlog we have 
    $$(h_{\alpha_x}, D(h_{\alpha_x})) \le (h_{\alpha_y}, D(h_{\alpha_y}))$$
    Hence $x, y \in D(h_{\alpha_y})$ subspace of $V$ and hence $x + \lambda y \in D(h_{\alpha_y}) \subseteq D$ so $D$ is a subspace of $V$. 
    Now we define $h \colon D |to \R$ as follows, for $x \in D$, then $x \in D(h_{\alpha_x})$ for some $\alpha_x \in I$, we define 
    $$h(x) = h_{\alpha_x}  (x)$$
    \begin{itemize}
        \item Well defined: If $x \in D(h_{\alpha_1}) \cap D(h_{\alpha_2})$ wlog we have $(h_{\alpha_1}, D(h_{\alpha_1})) \le (h_{\alpha_2}, D(h_{\alpha_2}))$ which gives us that 
        $h_{\alpha_2}(x) = h_{\alpha_1}(x)$. 

        \item $h$ is linear: this is because each $h_\alpha$ is linear.
        \item $h$ is bounded: Take $x\in D$ then $x \in D(h_{\alpha_x})$ so 
        $$\abs{h(x)} \le \norm{h_{\alpha_x}} \norm{x} \le \norm{f} \norm{x}$$
        so $h$ is 
    \end{itemize}
    todo, Zorn's lemma stuff.
\end{proof}

\begin{lemma}
    $V$ vector space over $\C$ and $f \colon V \to \C$ is linear $\iff f(x) = f_1(x) - if_1(ix)$ with $f_1 \colon V \to \R$ $\R$-linear. Moreover, $f$ is bounded iff $f_1$ is bounded with 
    $$\norm{f}_* = \norm{f_1}_*$$ 
\end{lemma}

\begin{proof}
    \begin{itemize}
        \item[$(\implies)$] $f \colon V \to \C$, then we can write $f(x) = f_1(x) + if_2(x)$ where $f_1, f_2 \colon V \to \R$. Therefore 
        $$f_1(x) = \frac{f(x) + \overline{f(x)}}{2}$$
        We want to show that $f_1$ is $\R$-linear, so take $\alpha \in \R$ and $x_1, x_2 \in V$ then 
        \begin{align*}
            f_1(x + \alpha y) &= \frac{f(x+ \alpha y) + \overline{f(x + \alpha y)}}{2} \\
            &= \frac{f(x) + \alpha f(y) + + \overline{f(x)} + \overline{\alpha} \overline{f(y)}}{2} \\
            &= \frac{f(x) + \overline{f(x)}}{2} + \alpha \frac{f(y) + \overline{f(y)}}{2} \\
            &= f_1 (x) + f_1(y)
        \end{align*}
        Moreover since $f$ is linear we have that 
        \begin{align*}
            f(ix) &= i f(x) \\
            f_1(ix) + if_2(x) &= if_1(x) - f_2(x) 
        \end{align*}
        Therefore $f_2(x) = if_1(ix)$, so 
        $$f(x) = f_1(x)  - i f_1(ix)$$

        \item[$\impliedby$] Easy.
    \end{itemize}

    For the moreover: 
    \begin{itemize}
        \item[$\implies$] $\abs{f_1(x)} \le \abs{f(x)} \le \norm{f}_* \norm{x}$ therefore $f_1$ is bounded with $\norm{f_1}_* \le \norm{f}_*$. Now 
        take $x \in V$, so $\exists \Theta_x$ s.t. $f(x) = \abs{f(x)} e^{i \Theta_x}$, therefore 
        \begin{align*}
            \abs{f(x)} &= f(x) e^{- i \Theta_x} \\
            &= f(e^{- i \Theta_x} x) \\
            &= f_1(e^{- i \Theta_x} x) + i f_2 (e^{i \Theta_x} x)  
        \end{align*}
        Comparing real and imaginary parts, we get 
        \begin{align*}
            \abs{f(x)} = f_1(e^{- i \Theta_x} x) \\
            &\le \norm{f_1}_* \norm{x}
        \end{align*}
        So 
        $$\norm{f} \le \norm{f_1}$$

        \item[$\impliedby$] easy. 
    \end{itemize}
\end{proof}

\begin{theorem}[Hahn-Banach] 
    $V$ normed space over $\F \in \set{\R, \C}$ and $M$ subspace of $V$ and $f \colon M \to \F$ bounded linear functional. Then $\exists F \colon V \to \F$ bounded linear functional 
    s.t. $F_{|M} = f$ and $\norm{F}_{V^*} = \norm{f}_{M^*}$.
\end{theorem}

\begin{proof}
    If $\F = \R$, done. If $\F = \C$, write $f(x) = f_1(x) - if_1(ix)$ with $f_1 \colon M \to \R$, so we can extend to $F_1 \colon V \to \R$ with 
    $$\norm{F_1} = \norm{f_1} =\norm{f}$$
    Define 
    $$F(x) = F_1(x) - iF_1(ix)$$ 
\end{proof}

\begin{corollary}
    Take $x_0 \in V$ with $x_0 \ne 0$, then $\exists f \in V^*$ s.t. $\norm{f} = 1$ and $f (x_0) = \norm{x_0}$.
\end{corollary}

\begin{proof}
    Take $M = \Span \set{x_0}$ with $g \colon M \to \F$ defined by 
    $$g(tx_0) = t \norm{x_0}$$
    Now take $x \in M$ so $x = tx_0$ then 
    $$\abs{g(x)} = \abs{t \norm{x_0}} = \norm{tx_0} = \norm{x}$$
    Therefore 
    $$\norm{g}_* = 1$$
    Apply Hahn-Banach.
\end{proof}

\textbf{Consequence:} If $V \ne \set{0}$, then $V^* \ne \set{0}$. 

\begin{corollary}
    $V \ne \set{0}$ normed space, and $x_0 \in V$. Then 
    $$\norm{x_0} = \sup_{\substack{f \in V^* \\ \norm{f}_* 1}} \abs{f(x_0)} = \max_{\substack{f \in V^* \\ \norm{f}_* = 1}} \abs{f(x_0)}$$
\end{corollary}

\begin{proof}
    Let $f \in V^*$ s.t. $\norm{f}_* = 1$, then 
    $$\abs{f(x_0)} \le \norm{x_0} \implies \sup_{\substack{f \in V^* \\ \norm{f}_* = 1}} \abs{f(x_0)} \le \norm{x_0}$$
    But by the corollary above we have that we can attain this sup. 
\end{proof}

\begin{corollary}
    $x \ne y$ in $V$, then $\exists f \in V^*$ s.t. $f(x) \ne f(y)$.
\end{corollary}

\begin{proof}
    $x \ne y$, so $x-y \ne 0$ in $V$. So by the above, $\exists f \in V^*$ s.t. 
    $$f(x-y) = \norm{x-y} = f(x) - f(y)$$
    as desired.
\end{proof}

\begin{corollary}
    $V$ vector space over $\F$ and $M$ proper closed subspace of $V$, and $x_0 \not\in M$ then $\exists f \in V^*$ s.t. $f_{|M} = 0$ and $f(x_0) = 1$.
\end{corollary}

\begin{proof}
    Let $g \colon M \oplus \Span \set{x_0} \to \F$ s.t. 
    $$g(m + tx_0) = t$$
    Then $g(x_0) = 1$ and clearly linear. Moreover take $x = m + tx_0$, with $\norm{x} = 1$, then for $t \ne 0$ we have 
    \begin{align*}
        \norm{x} &= 1 \\
        \abs{t} \norm{m/t + x_0} &= 1 \\
        \abs{t} &\le \frac{1}{d(x_0, M)}
    \end{align*} 
    for $t = 0$ we have $g(m) = 0$, so 
    $$\abs{g(x)} \le \frac{1}{d(x_), M}, \quad \forall x \in M \oplus \Span \set{x_0}$$
    Then $g$ is bounded, and $g_{|M} = 0$ and $g(x_0) = 1$. We Can extend by Hahn-Banach to some $f \in V^*$. 
\end{proof}

\begin{example}
    Let $x \in V$ with $f(x) = 0, \ \forall f \in V^*$. Then $x=0$. This is because 
    $$\norm{x} = \max_{\substack{f \in V^* \\ \norm{f} = 1}} \abs{f(x)}$$

    Now take $\varphi \colon (\ell^1) \to (\ell^\infty)^*$, so 
    $$\varphi(x) \colon b \to \sum_{n=1}^\infty b_n x_n$$
    with $b \in \ell^\infty$, then $\varphi$ is not surjective. This is because $C_0$ is a closed subspace of $\ell^\infty$, and $1 = (1,1,\cdots) \in \ell^\infty \setminus C_0$, 
    then $\exists T \in (\ell^\infty)^*$ s.t. $T_{|C_0} = 0$ and $T(1)= 1$. Suppose that $\varphi$ is surjective, then $\exists x \in \ell^1$ s.t. $\varphi(x) = T$. 
    So 
    $$Tb = \sum_{n=1}^\infty b_n x_n, \quad \forall b \in \ell^\infty$$
    Then since $e_i \in C_0$
    $$Te_i = x_i = 0$$
    So $x = 0$, hence $T = 0$, contradiction.
\end{example}

\begin{example}
    $V, W$ normed space and $V \ne \set{0}$, then $B(V, W)$ is Banach $\iff$ $W$ is Banach.    

    We already know $\impliedby$. Let $w_n$ be a Cauchy sequence in $W$. Take $v_0 \in V$ s.t. $\norm{v_0} = 1$, then $\exists f \in V^*$ with $\norm{f}_* = 1$ and 
    $f(v_0) = 1$. Now $w \in W$, define $T_w \colon V \to W$ given by 
    $$T_w (x) = f(x) \cdot w$$
    Now $T_w$ is linear, and 
    $$\norm{T_w x}_W = \abs{f(x)} \norm{w} \le \norm{x}_V \norm{w}_W$$
    Since this holds for all $x \in V$, we have $T_w \in B(V,W)$ and 
    $$\norm{T_w} \le \norm{w}$$
    Moreover, 
    \begin{align*}
        \norm{T_w v_0} &= \abs{f(x_0)} \norm{w} \\
        &= \norm{w}
    \end{align*}
    with $\norm{v_0} = 1$, so 
    $$\norm{T_w} = \norm{w}$$
    i.e $w \to T_w$ is an isometry. Now since $w_n$ is Cauchy in $W$, then $Tw_n$ is Cauchy in $B(V,W)$, so $Tw_n$ converges to some $T$ in $B(V,W)$. 
    Take $w = Tv_0$, now notice that 
    \begin{align*}
        \norm{w_n - w} &= \norm{f(v_0) w_n - f(v_0) w} \\
        &= \norm{T_{w_n} (v_0) - T (v_0)} \\
        &\le \norm{T_{w_n} - T} 
    \end{align*}
    Therefore $w_n \to w$.
\end{example}

\begin{proposition}
    $V$ normed space, if $V^*$ is separable, then $V$ is separable.
\end{proposition}

\begin{remark}
    Note that the converse isn't true since $\ell^1$ is separable and $(\ell^1)^* = \ell^\infty$ is not.
\end{remark}

\begin{corollary}
    $(\ell^\infty)^*$ is not isometrically isomorphic to $\ell^1$.
\end{corollary}

\begin{proof}
    $V^*$ is separable, then $V^*$ has a countable dense subset $\set{f_n}$. Now 
    \begin{align*}
        \norm{f_n} &= \sup_{\substack{x \in V \\ \norm{x} = 1}} \abs{f_n (x)}
    \end{align*}
    So $\exists x_n \in V$ with $\norm{x_n} = 1$ and 
    $$\frac{\norm{f_n}_*}{2} \le \abs{f_n(x_n)} \le \norm{f_n}_*$$
    Let $M = \Span \set{x_n}$, we claim that $\overline{M} = V$. 
    Suppose not, then $\exists x_0 \in V$ and $x_0 \not\in \overline{M}$, therefore $\exists f \in V^*$ with $f(x_0) = 1$ and $f_{|\overline{M}} = 0$,
    therefore $f(x_n) = 0$ $\forall n$. We want to show that this implies that $f = 0$. Let $\varepsilon > 0$, and $\exists N$ s.t. 
    \begin{align*}
        \norm{f_N - f} < \varepsilon/3
    \end{align*}
    So for $\norm{x} = 1$, 
    \begin{align*}
        \abs{f(x)} &\le \abs{f(x) - f(x)} + \abs{f_N(x)} \\
        &\le \norm{f - f_N} + \norm{f_N} \\
        &\le \norm{f -f_N} + 2 \abs{f_N(x_N)} \\
        &= \norm{f -f_N} + 2 \abs{f_N(x_N) - f(x_N)} \\
        &\le \norm{f -f_N} + 2 \norm{f - f_N} \\
        &< \varepsilon
    \end{align*}
    Therefore 
    $$\abs{f(x)} = 0 \implies \norm{f} = 0$$
    Hence $f = 0$, contradiction.

    Now $\overline{\Span \set{x_n}} = V \implies \overline{\Span_\Q \set{x_n}} = V$ so $V$ is separable. 
\end{proof}

\subsection{Adjoint Linear Transformation} 

\begin{definition}
    $V_1, V_2$ normed space and $T \colon V_1 \to V_2$ a bounded linear transformation. We define the adjoint of $T$ to be a map 
    $$T^* \colon V_2^* \to V_1^*, \quad (T^* f)  (x) =  f \circ T$$ 

    Moreover, we can show that $T^*$ is bounded, with 
    $$\norm{T^*} = \norm{T}$$
\end{definition}

\begin{remark}
    In the moreover use Hahn-Banach.
\end{remark}

\begin{remark}
    We have that 
    $$(T^* f, x)_{V_1^*, V_1} = (f, Tx)_{V_2^*, V_2}$$
    With $T^* \in \mathcal{L}(V_2^*, V_1^*)$, moreover 
    $$T \in B(V_1, V_2) \iff T^* \in B(V_2^*, V_1^*), \quad \norm{T^*} = \norm{T}$$
    Now 
    $$(T + \alpha S)^* = T^* + \alpha S^*, \quad \alpha \in \F$$
    And 
    $$(T \circ S)^* = S^* \circ T^*$$ 
    In the homework we prove that $T$ is an isomorphism iff $T^*$ is an isomorphism. 
\end{remark}

\begin{example}
    The Right-Shift operator 
    $$R \colon \ell^1 \to \ell^1, \ (x_1, x_2, \cdots) \to Rx = (0, x_1, x_2, \cdots)$$
    With $\norm{R} = 1$, we want to find $R^*$. We know that 
    $$R^* \colon (\ell^1)^* \to (\ell^1)^*$$
    Take $f \in (\ell^1)^*$, so there exists $b \in \ell^\infty$ s.t. 
    $$f(x) = \sum_{i=1}^\infty b_n x_n := f_b(x)$$
    Now 
    \begin{align*}
        (R^*f, x)_{(\ell^1)^*, \ell^1} &= (f_b, Rx) \\
        &= b_2 x_1 + b_3 x_2 + \cdots \\
        &= \sum_{i=1}^\infty b_{n+1} x_n 
    \end{align*}
    Therefore 
    $$R^* f_b = f_{(b_2, b_3, \cdots)} = f_{L(b)}$$
    Abusing notation we can say that 
    $$R^* b = L b$$    
\end{example}

\subsection{Reflexive Spaces}

\begin{definition}
    Let $V$ a normed space and $V^*$ be the dual of $V$, $B(V, \F)$ and $V^{**}$ the dual of $V^*$ $B(V^*, \F)$. Define the canonical map (the evaluation map)
    $$J_V \colon V \to V^{**}, \ (J_V x)  (f) = f(x)$$
    i.e. 
    $$(J_1 x, f)_{V^{**}, V^*} = (f, x)_{V^*, V}$$
\end{definition}

\begin{itemize}
    \item $J_V$ is linear: 
    \begin{align*}
        (J_V(x + \alpha y), f) &= f(x + \alpha y) \\
        &= f(x) + \alpha f(y) \\
        &= (J_V x, f) + \alpha (J_V y, f)
        &= (J_V(x) + \alpha J_V (y), f)
    \end{align*}

    \item \begin{align*}
        \norm{J_V x}_{V^{**}} &= \sup_{\substack{f \in V^{*} \\ \norm{f}_* = 1}} \abs{(J_V x, f)} \\
        &= \sup_{\substack{f \in V^* \\ \norm{f}_* = 1}} \abs{f(x)} \\ 
        &= \norm{x}
    \end{align*}
    Where the last equality is by Hahn-Banach, so $J_V$ is an isometry. 
    Therefore 
    $$V \cong J_V(V)$$
    Therefore $V^**$ contains an isomorphic copy of $V$. 
\end{itemize}

\begin{definition}
    We say that $V$ is reflexive iff $J_V$ is surjective. 
\end{definition}

\begin{remark}
    All reflexive spaces are Banach.
\end{remark}

\begin{example}
    \begin{itemize}
        \item $\ell^1$ is not reflexive, since if it was reflexive, then $(\ell^1)^{**} \cong \ell^1$, but 
        $(\ell^1)^{**} \cong (\ell^\infty)^*$ but this would mean that $\ell^\infty$ is separable, contradiction. 

        \item $\ell^p$ is reflexive for $1 < p < \infty$. Let 
        $$J_p \colon \ell^p \to (\ell^p)^{**}$$
        We want to show that this is surjective with 
        \begin{align*}
            (J_p x, f) &= (f, x)\\
            &= \sum_{n=1}^\infty b_n x_n \\
            &= \sum_{n=1}^\infty x_n b_n \\
            &= (f_x, b)_{(\ell^q)^*, \ell^q} \\
            &= (J_q b, f_x)_{(\ell^q)^{**}, (\ell^q)^*} 
        \end{align*}
        Take $\xi \in (\ell^p)^{**}$, now 
        \begin{align*}
            \xi \colon (\ell^p)^* \to \R 
        \end{align*}
        and we have some isomorphism $\varphi \colon \ell^q \to (\ell^p)^*$ with $\varphi(b) = f_b \colon x \in \ell^p \to \sum b_n x_n$. 
        Now we have an isomorphism
        $$\varphi^* \colon (\ell^p)^{**} \to (\ell^q)^*$$
        Now 
        \begin{align*}
            \varphi^* (\xi) \in (\ell^q)^*
        \end{align*} 
        so $\exists x \in \ell^p$ s.t. $\varphi^*(\xi) \colon b \in \ell^q \to \sum b_n x_n$. Now we need to show that 
        $$J_p (x) = \xi$$
        Indeed, for $f \in (\ell^p)^*$ we have 
        \begin{align*}
            (\xi, f) &= (\xi, \varphi(b)) \\
            &= (\varphi^* \xi, b) \\
            &= \sum b_n x_n \\
            &= \sum x_n b_n \\
            &= f_b (x) \\
            &= (J_V x, f_b) \\
            &= (J_Vx , f)
        \end{align*}
    \end{itemize}
\end{example}


\begin{lemma}%Note: AI-Generated.
    Let $V_1,V_2$ be normed vector spaces and let
    $T:V_1\to V_2$ be a bounded linear operator.
    Then the following diagram commutes:
    \[
    \begin{tikzcd}
        V_1 \arrow[r, "T"] \arrow[d, "J_{V_1}"']
        & V_2 \arrow[d, "J_{V_2}"] \\
        V_1^{**} \arrow[r, "T^{**}"']
        & V_2^{**}
    \end{tikzcd}
    \]
    That is,
    \[
        T^{**}\circ J_{V_1}=J_{V_2}\circ T.
    \]

    Moreover, if $T$ is a bounded linear isomorphism
    with bounded inverse, then $T^*$ is also an
    isomorphism and
    \[
        (T^*)^{-1}=(T^{-1})^*.
    \]
\end{lemma}


\begin{corollary}
    If $V_1 \cong V_2$, then $V_1$ is reflexive $\iff$ $V_2$ is reflexive. 
\end{corollary}

\end{document}
